% =============================================================================
% Project : THE_PHILOSOPHY_OF_VIABILITY
% Book    : Book I — Ontology of Viability
% Book ID : VS.B01
%
% File    : 2_MAINMATTER/PART_I_IDENTITY_AND_MANIFESTATION/CHAPTER_02_INVARIANCE/00_chapter_02.tex
%
% Chapter : 02 — Invariance
%
% Purpose :
%   Introduces the ontological problem of invariance, establishes
%   the scope and limits of the chapter, assembles its eleven
%   sections, and synthesizes the conclusions without reducing
%   invariance to identity, persistence, continuity, or immutability.
%
% Conceptual Scope :
%   Invariance, transformation, change, identity, difference,
%   structure, composition, continuity, persistence, and the
%   limits of ontological inference.
%
% =============================================================================

\chapter{Invariance}
\label{ch:invariance}

% -----------------------------------------------------------------------------
% CHAPTER METADATA
% -----------------------------------------------------------------------------

\newcommand{\ChapterTwoTitle}{Invariance}
\newcommand{\ChapterTwoID}{VS.B01.P01.C02}

% -----------------------------------------------------------------------------
% OPENING STATEMENT
% -----------------------------------------------------------------------------

\begin{quote}
Invariance concerns what remains unchanged under specified
transformations and in a specified respect; it does not, by
itself, determine the identity or persistence of that to
which the invariant aspect belongs.
\end{quote}

% =============================================================================
% CHAPTER INTRODUCTION
% =============================================================================

\section*{Introduction}
\addcontentsline{toc}{section}{Introduction}

The preceding chapter established the foundational problem
of identity and examined its relation to being. The present
chapter addresses invariance as a distinct ontological problem
that may intersect with identity but cannot be treated as
identical to it.

The problem of invariance arises when an existent, or some
determinate aspect of an existent, is considered in relation
to change or transformation. Under a specified transformation,
some aspect may remain unchanged while other aspects alter.
The identification of such an unchanged aspect gives rise
to a question about the conditions under which invariance
can meaningfully be asserted.

The central question of this chapter is therefore not simply
whether something changes or remains unchanged. It is:
what is claimed to remain unchanged, under which transformations,
within which domain, and in what respect?

For the purposes of the present inquiry, invariance is
approached through a relational formulation. An invariance
claim concerns an existent, a determination, or another
specified subject of assessment in relation to a defined
range of transformations, with respect to which a particular
aspect is claimed to remain unchanged.

This formulation identifies several analytically distinguishable
elements:

\begin{enumerate}
    \item \textbf{Subject:} the existent, determination, relation,
    structure, or other specified item to which the invariance
    claim refers.

    \item \textbf{Transformation:} the change, operation, transition,
    or class of changes under consideration.

    \item \textbf{Domain:} the defined range or field within which
    the transformations and the invariance claim are assessed.

    \item \textbf{Respect:} the particular feature, determination,
    relation, or condition with respect to which something is
    claimed to remain unchanged.

    \item \textbf{Criterion:} the basis on which the relevant
    sameness or preservation is assessed.
\end{enumerate}

These elements clarify the conditions of an invariance claim.
Their analytical distinction does not imply that they are
independently existing ontological entities. Nor does the
relational formulation, by itself, establish a complete
metaphysical account of the ground or nature of invariance.

The scope of an invariance claim must be proportionate to
its specified conditions. A feature that remains unchanged
under one transformation need not remain unchanged under
another. An aspect preserved within one domain need not be
preserved across a different domain. Likewise, a claim of
invariance with respect to one determination does not establish
the invariance of the subject in every respect.

Invariance must consequently be distinguished from absolute
immutability. The fact that something remains unchanged in
a specified respect does not entail that it undergoes no
change whatsoever. Conversely, the alteration of one aspect
does not establish that every other aspect has also changed.

The chapter also distinguishes invariance from identity,
self-sameness, persistence, continuity, stability, regularity,
symmetry, and conservation. These concepts may be relevant
to particular invariance claims, but none can be substituted
for invariance without further argument. In particular,
the recognition of an invariant feature does not by itself
establish numerical identity, and the absence of an identified
invariant does not establish the absence of identity.

The distinction between formal and ontological invariance
is equally important. A formal system may specify an invariant
under a set of operations according to its own rules and
criteria. Such a result establishes a formal relation within
that system. It does not automatically establish that the
same invariant has an ontological status independent of
the formal framework. Any transition from formal invariance
to an ontological claim requires additional justification.

The present chapter does not presuppose that invariance
must be grounded in an immutable substance, an underlying
essence, a fixed structure, or a single universal principle.
Nor does it assume that every existent must possess an
invariant aspect identifiable under every possible
transformation. Such claims would exceed what the relational
analysis alone can establish.

The chapter proceeds by examining the ontological problem
of invariance, its relation to identity, the conditions
required for defining invariants, and its connections with
transformation, difference, structure, composition, and
continuity. It then considers forms and domains of invariance,
the limits of invariance claims, and the conclusions that
can be retained without extending them beyond their grounds.

The purpose is to establish a disciplined conceptual account
of invariance and its inferential limits. The analysis is
not intended to settle every question concerning the ultimate
ontological status of invariance. It establishes distinctions
and conditions that must be respected in any further account
of its significance within the Ontology of Viability.

% =============================================================================
% SECTION 01 — ONTOLOGICAL PROBLEM OF INVARIANCE
% =============================================================================


% =============================================================================
% Section 01 — The Ontological Problem of Invariance
% Book I — Ontology of Viability
% Chapter 02 — Invariance
%
% Purpose:
%   To establish invariance as an independent ontological problem
%   and distinguish it from identity, persistence, continuity,
%   stability, and immutability.
%
% Conceptual Policy:
%   1. Invariance is not a definition or criterion of identity.
%   2. Invariance is not equivalent to the absence of change.
%   3. The existence of an invariant is not universally presupposed.
%   4. Every claim of invariance requires specification of its
%      domain and respect.
%   5. Ontological invariance is not reducible to epistemic
%      recognition, measurement, or description.
% =============================================================================

\section{The Ontological Problem of Invariance}
\label{sec:ontological-problem-of-invariance}

\subsection{The Problem of What Remains}

Existence is not necessarily characterized by an absence of change.
An existent may undergo transformations in its determinations,
configuration, composition, relations, or manifestations while
remaining an existent.

The occurrence of change, however, does not by itself establish
that everything associated with an existent changes. Nor does
the persistence of some determination or relation establish
that the existent as a whole remains unchanged.

This raises a distinct ontological problem:

\begin{quote}
What, if anything, remains invariant in an existent when it
undergoes transformation, and under what conditions can such
invariance be said to obtain?
\end{quote}

The problem concerns neither change alone nor the mere persistence
of an existent. It concerns the possible preservation of a
determination, relation, configuration, or other ontologically
relevant aspect across a specified transformation.

Invariance therefore requires an analysis of the relation between
what undergoes transformation and what, within the relevant
scope of that transformation, remains preserved.

The question is not whether an existent is simply the same
before and after a change. Such a formulation would risk
presupposing an answer to the problem of identity before
the relevant ontological conditions have been established.

Rather, the question concerns the possible preservation of
particular aspects of an existent under transformation,
without assuming that such preservation is equivalent to
the persistence of the existent's identity.

\subsection{Invariance as an Independent Ontological Problem}

Invariance is often discussed in connection with mathematical
structures, physical laws, symmetries, conservation principles,
and the behavior of systems under transformation.

These contexts provide important analytical resources.
They demonstrate that different properties, relations, or
structural features may remain preserved under specified
transformations.

However, the existence of such formal or scientific descriptions
does not, by itself, settle the ontological status of invariance.

A mathematical invariant is defined relative to a mathematical
structure and a specified class of transformations.
A physical invariant is established within a physical theory
and its relevant conditions.
A structural invariant may concern the preservation of an
organizational pattern despite changes in its components.

These examples indicate different possible forms of preservation,
but they do not establish that all such forms possess the same
ontological status or can be reduced to a single underlying
mechanism.

The ontological problem therefore requires a distinction
between the formal identification of an invariant and the
ontological interpretation of what is preserved.

An invariant cannot be regarded as ontologically established
merely because a quantity, relation, or expression remains
formally unchanged within a particular representation.

The relevant question is whether the preservation corresponds
to an actual determination or relation of the existent,
rather than exclusively to a feature of the formal framework
used to describe it.

This distinction does not diminish the significance of
mathematical or scientific results. It identifies the further
ontological question that must be addressed before those
results can support claims about the nature of existence.

\subsection{The Relational Character of Invariance}

Invariance cannot be adequately understood as an isolated
property attributed to an existent without reference to
the conditions under which preservation is asserted.

A claim of invariance necessarily involves several analytically
distinguishable elements:

\begin{enumerate}
    \item \textbf{The subject of invariance:}
    The existent, determination, relation, or configuration
    whose possible preservation is under consideration.

    \item \textbf{The transformation:}
    The change or transformation against which preservation
    is assessed.

    \item \textbf{The domain of transformation:}
    The range of conditions, states, configurations, or
    circumstances within which the transformation is considered.

    \item \textbf{The respect of invariance:}
    The particular determination, relation, magnitude,
    or organizational feature that is claimed to remain
    preserved.
\end{enumerate}

These elements are analytically distinguishable but need not
correspond to four independently existing entities.

In particular, the domain of transformation and the respect
of invariance specify the scope and terms of the claim.
They do not necessarily constitute additional ontological
components of the existent.

The relational character of invariance can be expressed
provisionally as follows:

\begin{quote}
An existent or one of its determinations is invariant
with respect to a specified transformation when the
relevant determination or relation remains preserved
throughout the transformation within the specified domain
and under the specified respect.
\end{quote}

This formulation is provisional because the ontological
meaning of preservation, the admissibility of transformations,
and the status of the relevant domain require further analysis
in the subsequent sections.

It also leaves open whether invariance should ultimately
be understood as a relation, a condition of preservation,
a feature of an existent, or a more complex ontological
configuration.

The relational formulation is therefore an analytical
starting point rather than a final ontological definition.

\subsection{Invariance and the Possibility of Change}

The existence of change does not necessarily exclude
invariance.

An existent may undergo changes in some determinations
while preserving others. Its composition may change while
a specified structural relation remains preserved.
Its quantitative magnitudes may vary while a particular
qualitative determination remains unchanged.
Its components may be replaced while a certain organizational
pattern continues to obtain.

These possibilities demonstrate that change and invariance
are not necessarily contradictory.

Their compatibility, however, depends on the respect under
consideration.

A determination may be invariant under one transformation
and non-invariant under another. A relation may be preserved
within one domain but cease to obtain beyond its boundaries.
A configuration may remain invariant under certain changes
while being altered by others.

Consequently, an assertion of invariance without specifying
the relevant transformation and domain is ontologically
incomplete.

The same applies to the opposite assertion. The fact that
an existent undergoes transformation does not justify
the conclusion that every determination, relation, or
structural feature associated with it has changed.

Invariance and change must therefore be examined in relation
to one another without reducing either to the other.

The relation between an initial configuration, a transformation,
and a resulting configuration can be represented schematically.
The representation does not presuppose that an invariant
must exist. It indicates that preservation and variation
are possibilities to be assessed with respect to specified
determinations and within a specified scope.


The figure illustrates that a transformation may be associated
with both possible preservation and possible variation.
Neither outcome should be assumed in advance. Whether a
determination remains invariant must be established through
an analysis of the relevant transformation, domain, and
respect of preservation.

\subsection{Distinguishing Invariance from Related Ontological Concepts}

The ontological problem of invariance requires distinguishing
it from several concepts that are frequently associated
with preservation and persistence.

These concepts may overlap in particular cases, but they
are not interchangeable.

\paragraph{Invariance and Identity.}

Identity concerns the ontological identity of an existent.
Invariance concerns the possible preservation of a specified
determination or relation under transformation.

The preservation of an invariant does not, by itself,
establish the numerical identity of an existent.
Likewise, the absence of an identified invariant does not,
by itself, establish the absence of identity.

Invariance therefore cannot serve as a general definition
or sufficient criterion of identity.

\paragraph{Invariance and Persistence.}

Persistence concerns the continued existence of an existent
through some interval or succession of conditions.

Invariance concerns what remains preserved under a specified
transformation.

An existent may persist while undergoing changes in many
of its determinations. Conversely, the preservation of
a particular invariant does not independently establish
the persistence of the existent as a whole.

\paragraph{Invariance and Continuity.}

Continuity concerns the manner in which existence, determination,
or transformation proceeds without a relevant rupture or
discontinuity.

Invariance concerns preservation under transformation.

A continuous process may involve extensive change, while
a discontinuous transformation may preserve certain
determinations or relations.

The two concepts therefore address different ontological
questions.

\paragraph{Invariance and Stability.}

Stability concerns the maintenance or boundedness of a
condition, configuration, or organization under relevant
conditions or disturbances.

Invariance concerns the preservation of a specified feature
under a specified transformation.

Stability may involve resistance to change, restoration,
or maintenance within a range of variation.
Invariance need not imply resistance, restoration, or
dynamic equilibrium.

\paragraph{Invariance and Immutability.}

Immutability denotes the absence of change in the relevant
respect.

Invariance does not necessarily require the absence of
change in the existent as a whole. It is compatible with
transformation, provided that the specified determination
or relation remains preserved.

Thus, invariance must not be reduced to immutability.

These distinctions establish that invariance occupies
a distinct analytical position within ontology.
Its relation to each of these concepts must be investigated
without assuming conceptual equivalence.

The distinction between invariance and its neighboring
ontological concepts does not imply that invariance itself
has only one form. Different determinations and relations
may be preserved under transformation, giving rise to
different analytical forms of invariance.

The following classification identifies several such forms.
It is an analytical classification rather than an assertion
that these forms are mutually exclusive, exhaustive, or
ontologically independent.


The forms presented in the table identify different respects
in which preservation may be examined. Their classification
does not establish that every existent exhibits all these
forms, nor that any particular form must obtain under every
transformation.

\subsection{The Scope and Limits of an Invariance Claim}

An assertion that something is invariant is incomplete
unless its scope is specified.

At minimum, such an assertion requires clarification of:

\begin{itemize}
    \item what is claimed to remain preserved;
    \item which transformation is relevant;
    \item which domain of transformation is considered;
    \item under which respect preservation is asserted;
    \item whether preservation is absolute within that domain
    or conditional upon particular circumstances.
\end{itemize}

Without these specifications, invariance risks becoming
an unrestricted attribution of permanence.

No universal assumption should be made that every existent
possesses an invariant, that every transformation preserves
some feature, or that every apparent preservation has
ontological significance.

An invariant may exist only under restricted conditions.
It may be preserved within one domain and cease to be
preserved beyond it.
Different determinations of the same existent may exhibit
different patterns of preservation.

Furthermore, the identification of an invariant within
a limited domain does not justify its generalization
to all possible transformations.

The scope of an invariance claim must therefore remain
proportionate to the conditions under which preservation
is established.

\subsection{The Ontological Problem and Its Subsequent Development}

The preceding analysis establishes invariance as a distinct
ontological problem concerning the possible preservation
of determinations and relations under transformation.

It does not yet establish a complete definition of invariance,
a final classification of its forms, or a universal account
of the conditions under which invariance obtains.

Those questions require separate examination.

The subsequent sections will develop the problem through
the following sequence:

\begin{enumerate}
    \item The relation between invariance and identity,
    with particular attention to their ontological distinction.

    \item The nature and defining conditions of invariants,
    including the criteria by which preservation may be
    distinguished from formal or apparent sameness.

    \item The relation between invariance and transformation,
    including the conditions under which change and
    preservation can coexist.

    \item The relation between invariance and difference,
    with attention to the possibility of preserved features
    within configurations that undergo differentiation.

    \item The role of structure and composition in the
    preservation of determinations and relations.

    \item The relation between invariance and continuity,
    without reducing either concept to the other.

    \item The forms and domains within which invariance
    may obtain.

    \item The limits of invariance and the conditions
    under which claims of preservation cease to be valid.
\end{enumerate}

The purpose of this development is not to establish
invariance as the ultimate foundation of identity or
existence.

Rather, it is to determine the ontological status,
scope, conditions, and limitations of preservation
within a world in which transformation is possible.

The central question remains:

\begin{quote}
What can remain invariant through transformation,
and what does such preservation mean for the ontological
constitution of the existent?
\end{quote}


% =============================================================================
% SECTION 02 — INVARIANCE AND IDENTITY
% =============================================================================

% =============================================================================
% Project : THE_PHILOSOPHY_OF_VIABILITY
% Book    : Book I — Ontology of Viability
% Book ID : VS.B01
%
% File    : 2_MAINMATTER/PART_I_IDENTITY_AND_MANIFESTATION/CHAPTER_02_INVARIANCE/sections/section_02_invariance_and_identity.tex
%
% Section / Chapter:
%   Section 02 — Invariance and Identity
%   Chapter 02 — Invariance
%
% Purpose :
%   To establish the ontological distinction between identity and
%   invariance, examine the limits of invariant preservation as a
%   criterion of identity, and clarify the possible relevance of
%   invariance to the characterization of an existent.
%
% Conceptual Policy:
%   1. Being and identity are ontologically inseparable but
%      analytically distinguishable.
%   2. Neither being nor identity has ontological priority over
%      the other.
%   3. Invariance is neither the definition nor the foundation
%      of identity.
%   4. The preservation of an invariant does not, by itself,
%      establish numerical identity.
%   5. The absence of an identified invariant does not establish
%      the absence of identity.
%   6. Multiple invariants do not entail multiple identities,
%      nor do they establish a single identity.
%   7. The possible plurality and relations of identities remain
%      open ontological questions.
%
% =============================================================================

\section{Invariance and Identity}
\label{sec:invariance-and-identity}

\subsection{The Need for an Ontological Distinction}

Invariance and identity become relevant to the same
ontological problem whenever an existent is considered
in relation to transformation. The analysis of
transformation may raise two questions: what, if
anything, remains preserved, and whether the existent
in question retains its identity.

These questions are related, but they are not identical.

Invariance concerns the preservation of a specified
determination, relation, or configuration under a
specified transformation and within a specified domain.
Identity concerns the ontological identity of the
existent. The former examines a possible relation
between a feature and a transformation; the latter
concerns the existent in its identity.

The distinction is not merely terminological. Each
question requires a different kind of ontological
clarification. Establishing that a determination
remains preserved does not, without further argument,
establish what the identity of the existent consists
in. Conversely, establishing that a determination
has changed does not, by itself, settle whether the
existent retains its identity.

A preserved feature may be relevant to an account
of an existent, but relevance must not be confused
with constitution or sufficiency. The fact that a
determination remains unchanged does not establish
that it constitutes the identity of the existent,
and the fact that several determinations are preserved
does not establish that their collection is identical
with that identity.

The ontological distinction must therefore be maintained
before any account of the relation between identity
and invariance can be developed. Without it, the
analysis risks replacing the problem of identity with
a narrower inquiry into the preservation of selected
features.

\subsection{The Ontological Status of Identity}

The preceding chapter established three distinct
principles concerning the relation between being
and identity: their analytical distinction, their
ontological inseparability, and the absence of
ontological priority between them.

These principles must remain distinct. Analytical
distinction permits being and identity to be examined
through different questions. Ontological inseparability
denies that they constitute independently existing
realities. The absence of ontological priority denies
that either can be established as ontologically
prior to the other.

Within this framework, identity is not an additional
property attached to an existent after its existence
has already been constituted. Nor is it a secondary
result produced by the preservation of particular
determinations.

An existent is not first ontologically complete
and subsequently supplied with identity through
the maintenance of an invariant configuration.
Such an account would make identity dependent upon
a process of preservation and would thereby reduce
an ontological question to a condition concerning
the behaviour of selected features.

Identity must instead be considered in relation
to the existent itself, without presupposing that
its ontological status can be derived from a list
of preserved determinations.

This does not mean that identity has ontological
priority over being. The distinction between the
ontological status of an existent and the preservation
of its features does not introduce a hierarchy between
being and identity. It establishes only that invariance
and identity address different questions.

Invariance concerns whether a specified feature
is preserved under a transformation. It does not
constitute an independent ontological ground from
which the identity of the existent can be derived.

\subsection{The Insufficiency of Invariance as a Criterion of Identity}

To reduce identity to invariance would be to maintain
that the identity of an existent is determined by
the feature or features that remain preserved under
a relevant class of transformations.

Such a reduction requires an additional premise:
that the preserved features selected for analysis
are constitutive of, or sufficient to establish,
the identity of the existent.

That premise cannot be obtained from preservation
alone.

A feature may remain unchanged without being
constitutive of the identity of the existent. A
quantitative determination may remain constant
while the composition of an existent changes.
A structural relation may be preserved while
particular components are replaced. A qualitative
determination may remain unchanged while other
aspects of the existent undergo substantial
modification.

These cases do not establish that identity is
unaffected by every such change. They establish
a more limited point: preservation of a feature
does not, by itself, determine the ontological
significance of that feature for identity.

The selection of an invariant as an identity criterion
therefore requires justification beyond the fact
that the feature remains preserved. It must be
established why that particular determination
or relation should be regarded as relevant to
identity, and why its preservation should be
sufficient for the intended conclusion.

Without such justification, the reduction remains
unsupported.

A further difficulty arises from the plurality
of possible invariants. Different determinations
and relations may remain preserved under different
transformations, or within different domains of
the same transformation. A quantitative feature,
a structural relation, and a relational configuration
may each exhibit preservation under specified
conditions.

Their preservation does not establish that they
are equivalent in ontological significance. Nor
does it provide a rule for deciding whether one
of them, several of them, or their combination
should count as constitutive of identity.

The problem cannot be resolved merely by collecting
more invariants. An aggregate of preserved features
does not become an account of identity simply
because the number of preserved features increases.

The reduction of identity to invariance therefore
faces a problem of justification: preservation
establishes a fact about a specified feature under
specified conditions, but it does not independently
establish the ontological role of that feature
in the identity of an existent.

\subsection{The Preservation of an Invariant and Numerical Identity}

Numerical identity concerns whether two designations
or configurations refer to one and the same existent,
rather than to two distinct existents that may be
qualitatively similar, structurally corresponding,
or otherwise alike.

The preservation of an invariant does not, by itself,
resolve this question.

Two numerically distinct existents may share the
same quantitative determination, possess corresponding
structural relations, or exhibit equivalent patterns
of organization. A feature may therefore be preserved
or shared across distinct existents without making
them numerically identical.

The preservation of a feature establishes sameness
in a specified respect. Numerical identity is not
equivalent to sameness in every respect, nor can
it be inferred from the preservation of one selected
respect.

The same distinction applies when an existent is
considered across transformation. If a determination
remains preserved between an initial and a resulting
configuration, this establishes preservation with
respect to that determination. It does not, without
a further ontological argument, establish that
the configurations belong to one and the same
existent.

The relation between configurations and the identity
of an existent must therefore be examined independently
of the preservation claim.

This distinction does not make invariance irrelevant
to the investigation of numerical identity. An
invariant may contribute to an argument concerning
the relation between configurations. It may provide
a relevant condition, constraint, or indication
within a broader account. What it cannot do is
establish numerical identity merely by being
preserved.

The evidential or explanatory relevance of an
invariant must consequently be distinguished from
its sufficiency as a criterion of numerical identity.

\subsection{The Absence of an Identified Invariant and the Question of Identity}

The converse reduction is equally unwarranted.
If identity were defined through invariant
preservation, the failure to identify an invariant
would appear to undermine the identity of the
existent.

That conclusion does not follow.

The absence of an identified invariant may mean
that no relevant feature has been shown to remain
preserved within the specified domain and under
the transformation examined. It may also indicate
that the selected respect is not preserved, that
the transformation exceeds the range within which
a proposed invariant applies, or that the criteria
of preservation have not been adequately specified.

None of these possibilities establishes the absence
of identity.

The distinction between an invariant not being
identified and an invariant not existing must also
be maintained. Failure to establish a preservation
claim is not equivalent to establishing that no
feature could be preserved under any relevant
transformation.

Even where a particular determination is shown
to change, the ontological question of identity
remains open. Changes in composition, structure,
relations, or other determinations may be relevant
to that question, but their significance cannot
be settled merely by recording the absence of
preservation in one respect.

Whether an existent retains its identity through
a particular transformation, and under what
conditions this is possible, requires an independent
ontological analysis.

Accordingly, neither the identification nor the
non-identification of an invariant can be treated
as a conclusive determination of identity.

\subsection{Invariance as a Dimension of Ontological Characterization}

Rejecting the reduction of identity to invariance
does not render invariance ontologically insignificant.

The preservation of a determination or relation
may contribute to the characterization of an
existent. It may disclose which features remain
unchanged under a specified transformation and
which features undergo alteration. Such distinctions
can be relevant to the analysis of constitution,
composition, relational organization, and persistence.

For example, the preservation of a structural
relation may establish that a particular organizational
feature remains unchanged despite modifications
in the components participating in that relation.

The preservation of a quantitative determination
may establish that a specified magnitude remains
constant within a particular domain.

The preservation of a relation among components
may establish that a particular relational arrangement
continues to obtain despite changes in other
determinations.

In each case, the invariance claim concerns a
specified respect. Its ontological significance
depends upon the feature examined, the transformation
considered, the domain within which the claim
applies, and the criterion by which preservation
is established.

These findings may contribute to a more extensive
account of an existent. They may constrain certain
interpretations of its transformation or clarify
the respects in which its configuration is maintained
or altered.

However, characterization must not be confused
with ontological constitution. To characterize
an existent through one or more preserved features
is not necessarily to establish what constitutes
its identity.

The distinction is especially important when
invariance is used to discuss persistence. A
determination may persist, a relation may remain
preserved, or a configuration may exhibit continuity
under specified conditions. None of these descriptions
is automatically equivalent to the persistence
or continuity of the existent itself.

The relation between these concepts requires
separate examination. Invariance may contribute
to their analysis, but it cannot substitute for
their ontological clarification.

\subsection{Multiple Invariants and the Open Question of Multiple Identities}

A single existent may exhibit different forms
of invariance under different transformations
or within different domains.

Some determinations may remain preserved while
others change. A structural relation may remain
invariant under one transformation but not under
another. A relational configuration may be
preserved within a restricted domain while
undergoing alteration under a broader class
of transformations.

This plurality of invariants raises questions
concerning the relation between different respects
of preservation and the identity of the existent.

It does not, however, establish a plurality of
identities.

There is no justified inference from the existence
of several invariants to the existence of several
identities. Each invariant concerns a specified
feature under specified conditions. The number
of such features does not determine the number
of identities to which they may be related.

Nor can the totality of invariants automatically
be treated as a single identity. Their aggregation
would still leave unresolved why the selected
features should constitute identity and how their
different domains, transformations, and respects
should be related.

The possible plurality of identities is a separate
ontological problem. It cannot be settled by
counting preserved features, identifying different
forms of invariance, or assigning an identity
to every level of constitution.

Whether an existent may possess multiple identities,
whether identity is to be considered at different
ontological levels, and how any such identities
might be related remain open questions requiring
independent argument.

The analysis of invariance may help formulate
these questions more precisely. It does not
provide their answer.

\subsection{The Analytical Independence of Identity and Invariance}

The distinction developed in this section can
be summarized through the questions addressed
by each concept.

Identity concerns the ontological identity of
an existent. It asks how the existent is to be
understood in its identity, without presupposing
that this question can be answered through the
preservation of selected features.

Invariance concerns whether a specified determination,
relation, or configuration remains preserved under
a specified transformation and within a specified
domain.

The two questions are analytically independent
in the sense that neither can be answered merely
by substituting the other for it.

This independence does not imply that identity
and invariance are unrelated. An invariance claim
may contribute to the characterization of an
existent, while an account of identity may constrain
the interpretation of preservation across
configurations.

Their relation, however, requires an argument
that respects the distinct subject matter of
each inquiry.

This distinction also limits the direct transfer
of formal definitions of invariance into ontology.
A formal framework may define identity through
equivalence under a specified class of allowable
transformations. Such a definition can be valid
within the framework for which it is constructed.

Its formal validity does not, by itself, establish
that the same criterion determines the ontological
identity of an existent.

An ontological application would require independent
justification of the relevant transformations,
the domain of application, the respects of
preservation, and the reason why the resulting
equivalence relation should be interpreted as
identity rather than as a more limited form
of correspondence or sameness.

Without that justification, a formal invariance
criterion remains a framework-specific construction
and must not be elevated into a general ontological
principle.

The distinction between identity and invariance
is therefore not a denial of their possible
relation. It is a requirement that this relation
be established without reducing one problem
to the other.

\subsection{Conclusion: Invariance Without Reduction to Identity}

Identity and invariance address distinct ontological
questions.

Identity concerns the existent in its ontological
identity. Invariance concerns the possible preservation
of a specified determination, relation, or configuration
under a specified transformation and within
a specified domain.

The preservation of an invariant does not, by
itself, establish numerical identity. The absence
of an identified invariant does not establish
the absence of identity. Multiple invariants
do not entail multiple identities, and their
aggregation does not automatically constitute
a single identity.

Nevertheless, invariance remains relevant to
ontological analysis. It can characterize particular
respects of preservation, clarify the relation
between transformation and alteration, and
contribute to inquiries concerning constitution,
composition, and persistence.

Its contribution must remain proportionate to
what an invariance claim actually establishes.
A claim concerning the preservation of a feature
must not be extended into a conclusion about
the identity of the existent without an additional
argument.

The distinction established in this section
therefore preserves the ontological independence
of the problem of identity while allowing
invariance to retain its proper analytical
significance.

The next section examines the conditions under
which a determination or relation may properly
be described as invariant, including the role
of transformation, domain, respect, and the
criterion of preservation.

% =============================================================================
% END OF SECTION 02 OF CHAPTER 02
% =============================================================================

% =============================================================================
% SECTION 03 — DEFINING INVARIANTS
% =============================================================================

% =============================================================================
% Project : THE_PHILOSOPHY_OF_VIABILITY
% Book    : Book I — Ontology of Viability
% Book ID : VS.B01
%
% File    : 2_MAINMATTER/PART_I_IDENTITY_AND_MANIFESTATION/CHAPTER_02_INVARIANCE/sections/section_03_defining_invariants.tex
%
% Section / Chapter:
%   Section 03 — Defining Invariants
%   Chapter 02 — Invariance
%
% Purpose :
%   To establish the analytical requirements for formulating and
%   evaluating claims that a determination or relation is preserved
%   under a specified transformation, while distinguishing the
%   formal specification of invariance from its ontological grounds.
%
% Conceptual Policy:
%   1. An invariance claim must identify what is claimed to be
%      preserved, under which transformation, within what domain,
%      and according to which criterion.
%   2. The subject of an invariance claim must be distinguished
%      from the existent to which it belongs.
%   3. The transformation and the domain of its applicability
%      must be specified without presupposing universal validity.
%   4. The respect of preservation must be distinguished from
%      the subject being examined.
%   5. Exact preservation, approximate correspondence, and
%      resemblance must not be conflated.
%   6. The admissibility of a transformation is a constraint
%      on the scope of a claim, not an independent ontological
%      component of invariance.
%   7. Formal invariance does not automatically establish
%      ontological invariance.
%   8. Analytical conditions specify and evaluate a claim;
%      they do not create the feature claimed to be preserved.
%   9. No invariant is presumed to hold across unrestricted
%      transformations or domains.
%
% =============================================================================

\section{Defining Invariants}
\label{sec:defining-invariants}

\subsection{The Problem of Defining an Invariant}

The assertion that something remains unchanged is not,
by itself, an adequate definition of invariance.

Such an assertion remains indeterminate unless it specifies
what is claimed to remain preserved, under which
transformation, within what domain, and according to
which criterion.

A determination may remain preserved under one
transformation while undergoing alteration under
another. A relation may remain unchanged within
a restricted domain but fail to remain so beyond
that domain. A configuration may preserve one
organizational feature while its other determinations
or relations undergo modification.

Consequently, invariance cannot be treated as an
unrestricted property of an existent, nor can it
be established through the undifferentiated description
of two configurations as unchanged.

An invariance claim is a claim of preservation under
specified conditions. Its meaning and validity depend
upon the identification of the feature concerned,
the transformation considered, the domain within
which the claim applies, and the criterion by which
preservation is assessed.

The task of defining invariants is therefore not
to introduce a universal mechanism of preservation.
It is to establish the analytical requirements
through which particular claims of invariance
can be formulated, distinguished, and evaluated.

These requirements must be sufficiently precise
to prevent the concept from being extended beyond
what a particular claim actually establishes.

\subsection{The Subject of an Invariance Claim}

Every claim of invariance identifies something
whose preservation is being examined. This is
the subject of the invariance claim.

The subject may be a determination, relation,
structural feature, quantitative magnitude,
qualitative feature, or another identifiable
aspect associated with an existent.

It must be distinguished from the existent as
a whole.

An existent may undergo substantial transformation
while a particular determination or relation
remains preserved. Conversely, the preservation
of one feature does not establish that the
existent as a whole remains unchanged.

The subject of an invariance claim therefore
identifies the particular feature under examination
without implying that this feature exhausts
the constitution of the existent or determines
its identity.

The distinction is also important when the
subject is relational. A claim concerning the
preservation of a relation must identify the
relation under consideration and the relevant
terms or relata involved in it. The preservation
of a relational form does not necessarily imply
the preservation of every determination of
the relata, nor does it establish that the
relata themselves remain numerically identical.

The subject must be sufficiently determinate
for the preservation claim to be evaluated.
An unspecified assertion that the existent,
its structure, or its organization remains
unchanged leaves open which feature is actually
being claimed as invariant.

Thus, identifying the subject establishes
what the claim concerns; it does not yet
establish that preservation occurs.

\subsection{Transformation as the Relational Context of Invariance}

Invariance is intelligible only in relation
to a transformation.

A transformation provides the context within
which the preservation or alteration of a
specified determination or relation can be
examined. Without a transformation, there
is no specified relation of preservation
to assess.

The transformation must be characterized
sufficiently to distinguish the relevant
initial and resulting configurations and
to identify the kind of change under
consideration.

A transformation may involve changes in
determinations, relations, composition,
structure, or other aspects of an existent.
It may also be restricted to a particular
class of changes.

Not every difference between two configurations
establishes that one has been transformed
into the other. The relation between the
configurations must itself be specified
where the claim depends upon a transformation
connecting them.

Likewise, not every transformation preserves
a feature. Preservation is a possible result
to be established, not an outcome that follows
from the mere specification of a transformation.

The same feature may remain preserved under
one transformation and undergo alteration
under another. Accordingly, an invariant
cannot be identified without reference to
the transformation relative to which its
preservation is claimed.

Transformation is therefore not an additional
property of the invariant. It specifies
the relation under which a preservation
claim is made.

\subsection{The Domain of Applicability}

A transformation is not necessarily applicable
to every possible configuration or condition.
The domain of an invariance claim specifies
the range within which the relevant
transformation and preservation criterion
are being considered.

A domain may be restricted by the configurations
included, the conditions under which the
transformation is defined, or the range
of cases to which the claim is intended
to apply.

For example, a structural relation may
remain preserved across a specified class
of transformations while undergoing
alteration under transformations outside
that class.

The preservation established within the
restricted domain does not automatically
extend beyond it.

The domain must therefore be distinguished
from the transformation itself. The
transformation specifies the relevant
change; the domain specifies the range
within which that transformation and
the corresponding preservation claim
are considered applicable.

A domain is an analytical delimitation
of a claim. Its specification does not,
by itself, establish an independent
ontological realm or introduce a new
kind of existent.

Nor does a restricted domain imply that
the claim is arbitrary. A domain may be
grounded in the actual conditions under
which a transformation occurs or in the
specified range of configurations to
which the claim applies.

The validity of the claim nevertheless
remains bounded by that domain unless
further grounds justify its extension.

\subsection{The Respect of Preservation}

The respect of preservation specifies
the regard in which a determination
or relation is claimed to remain unchanged.

It must be distinguished from the subject
of the claim.

The subject identifies what is being
examined. The respect specifies the
aspect under which its preservation
is assessed.

A configuration, for example, may retain
a particular structural arrangement while
its quantitative determinations change.
A relation may preserve its formal
organization while the properties of
its relata undergo alteration.

In such cases, the claim of preservation
is meaningful only if the relevant
respect is identified.

The respect must also be distinguished
from the criterion of preservation.
Identifying a respect establishes the
aspect in which sameness is being
examined; the criterion specifies
what counts as preservation in that
respect.

Without this distinction, a claim may
move ambiguously between different
features or different meanings of
sameness.

The preservation of a quantitative
determination, for example, cannot
be substituted without argument for
the preservation of a structural
relation. Likewise, the preservation
of a relational arrangement does
not establish the preservation of
the material constitution of the
terms participating in it.

An invariance claim must therefore
remain limited to the respect actually
specified and established.

\subsection{The Criterion of Preservation}

Identifying a subject, transformation,
domain, and respect does not yet establish
that preservation has occurred.

A criterion of preservation is required
to determine what must remain the same
for the specified subject to qualify
as invariant under the transformation
considered.

The criterion must correspond to the
nature of the subject and the respect
under examination.

Where the subject is a quantitative
determination, the criterion may require
the preservation of a specified value.
Where the subject is a structural relation,
the criterion may concern the maintenance
of a specified relational organization.
Where the subject is a configuration,
the criterion must identify which features
of that configuration are required to
remain preserved.

These examples do not establish a universal
criterion. The appropriate criterion
depends upon the subject and the claim
being formulated.

Several distinctions are necessary.

First, resemblance is not sufficient
to establish invariance. Two features
may appear similar without satisfying
the criterion required for preservation.

Second, correspondence is not necessarily
preservation. Two configurations may
correspond according to a mapping or
comparison rule while differing in
the respect relevant to the claim.

Third, approximate agreement must be
distinguished from exact preservation.
Where a criterion permits a tolerance
or degree of variation, that allowance
must be explicitly specified. A claim
of approximate invariance must not
be presented as exact invariance.

Fourth, persistence must not be substituted
for invariance. The continued occurrence
of a feature does not, by itself, establish
that the feature remains unchanged in
the respect under examination.

The criterion of preservation therefore
provides the standard against which
the claim is evaluated. It does not
create the preservation it is intended
to assess.

\subsection{The Scope and Admissibility of Transformation}

An invariance claim must specify which
transformations are included within
its scope.

Not every conceivable transformation
is relevant to every claim. A transformation
must be admissible relative to the
subject, domain, and respect under
consideration.

Admissibility concerns whether a
transformation belongs to the specified
class of transformations against which
the preservation claim is being assessed.

For example, a feature may be preserved
under a defined class of structural
transformations. This does not establish
that it remains preserved under changes
that fall outside that class or alter
the conditions required for the feature
to be identified in the first place.

The admissibility of a transformation
must therefore be justified by the
scope of the claim rather than assumed
from the mere possibility of describing
a change between two configurations.

This requirement does not introduce
a separate ontological component of
invariance. It clarifies the limits
of the transformation class relevant
to the claim.

The scope of validity is determined
by the combined specification of the
subject, transformation, domain, respect,
and preservation criterion, together
with the admissibility conditions
governing the transformations included.

A claim established under one set of
conditions cannot be generalized
automatically to a broader set.

\subsection{Formal Invariance and Ontological Invariance}

A distinction must be maintained between
formal invariance and ontological invariance.

A formal system may define an invariant
through a rule, operation, equation,
mapping, or transformation that preserves
a specified mathematical or logical
feature.

Such a result establishes preservation
relative to the formal structure and
the conditions defined within that
system.

An ontological claim, however, concerns
whether the feature in question is
actually preserved in the relevant
determinations or relations of an
existent.

The validity of a formal result and
the ontological status of the feature
represented by that result are therefore
analytically distinguishable.

A formal representation may describe,
model, or express an aspect of an
ontological relation. It may also
provide a rigorous means of investigating
whether a proposed feature remains
preserved under specified conditions.

Nevertheless, the formal preservation
of a mathematical or logical expression
does not automatically establish that
the corresponding feature is an
ontological invariant of an existent.

An ontological application requires
an account of the relation between
the formal structure and the existent
or feature to which the formal result
is being applied.

This requirement does not imply that
formal methods are ontologically
irrelevant. It establishes that their
results must not be extended beyond
their justified scope.

Formal invariance may contribute to
the investigation of ontological
invariance, but the transition from
a formal result to an ontological
claim requires an additional argument.

\subsection{The Analytical Conditions of an Invariance Claim}

The preceding analysis distinguishes
the principal elements required to
formulate and evaluate an invariance
claim.

These elements do not all have the
same logical function.

The subject, transformation, domain,
respect, and criterion specify the
content of the preservation claim.
The admissibility of transformations
and the scope of validity delimit
the conditions under which the claim
is applicable. Ontological grounding
concerns whether the result established
within an analytical or formal
framework warrants an ontological
interpretation.

They must therefore not be treated
as independent ontological components
or as successive stages through which
an invariant necessarily comes into
existence.

Their relations are summarized in
the following table.

% =============================================================================
% Project : THE_PHILOSOPHY_OF_VIABILITY
% Book    : Book I — Ontology of Viability
% Book ID : VS.B01
%
% File    : 4_ASSETS/tables/VS.B01.P01.C02.S03.SS10.P001.TL01.conditions_of_an_invariance_claim.tex
%
% Section / Chapter / Asset:
%   Chapter 02 — Invariance
%   Section 03 — Defining Invariants
%   Table TL01 — Conditions of an Invariance Claim
%
% Asset ID:
%   VS.B01.P01.C02.S03.SS10.P001.TL01.conditions_of_an_invariance_claim
%
% Purpose :
%   To distinguish the principal analytical elements and validity
%   conditions required for formulating and evaluating an invariance
%   claim.
%
% Conceptual Policy:
%   1. The table presents analytical requirements, not ontological
%      components.
%   2. Specification, scope, and ontological grounding have distinct
%      logical functions.
%   3. The table does not presuppose that an invariant exists.
%   4. The requirements are not presented as a universal formal
%      algorithm.
%
% =============================================================================

\begin{table}[htbp]
    \centering

    \caption{
        Analytical Conditions for Formulating and Evaluating
        an Invariance Claim
    }
    \label{tab:conditions-of-an-invariance-claim}

    \vspace{0.5em}

    \small

    \renewcommand{\arraystretch}{1.45}

    \begin{tabular}{
        |p{0.23\textwidth}
        |p{0.65\textwidth}|
    }
        \hline

        \textbf{Analytical Element}
        &
        \textbf{Question and Function}
        \\ \hline

        \textbf{Subject}
        &
        What determination, relation, structure, or other
        identifiable feature is claimed to be preserved?
        \\ \hline

        \textbf{Transformation}
        &
        Under which specified change or class of changes
        is preservation being examined?
        \\ \hline

        \textbf{Domain}
        &
        Within which range of configurations, conditions,
        or cases does the transformation and the claim apply?
        \\ \hline

        \textbf{Respect}
        &
        In what particular regard is the subject claimed
        to remain preserved?
        \\ \hline

        \textbf{Preservation Criterion}
        &
        What standard determines whether the specified
        subject remains unchanged in the relevant respect?
        \\ \hline

        \textbf{Transformation Admissibility}
        &
        Which transformations legitimately belong to the
        specified class against which preservation is assessed?
        \\ \hline

        \textbf{Scope of Validity}
        &
        What are the limits within which the preservation
        claim is justified, and beyond which it cannot
        be generalized without further argument?
        \\ \hline

        \textbf{Ontological Grounding}
        &
        What justifies interpreting the established
        preservation as a feature of the relevant existent
        or relation rather than merely as a formal result?
        \\ \hline

    \end{tabular}

    \normalsize

    \vspace{0.4em}

    \begin{minipage}{0.92\textwidth}
        \footnotesize
        \textit{Note.}
        The elements listed above have different analytical
        functions. They specify the content, applicability,
        limits, and ontological interpretation of an
        invariance claim. They are not independent
        ontological components and do not imply that
        an invariant necessarily exists.
    \end{minipage}

\end{table}

% =============================================================================
% END OF TABLE TL01
% =============================================================================

The table is an analytical guide
to the formulation of invariance
claims. It does not constitute an
exhaustive taxonomy of all possible
forms of invariance, nor does it
presuppose that every claim must
be expressed through an identical
formal procedure.

The requirements must be interpreted
in relation to the particular subject
and transformation under consideration.

A claim that identifies a subject
but leaves its transformation
unspecified remains incomplete.
A claim that specifies a transformation
but leaves its domain or respect
indeterminate does not establish
the full scope of preservation.
A claim that identifies a formal
invariant without justifying its
ontological application does not,
by that fact alone, establish an
ontological invariant.

The framework clarifies the conditions
of a meaningful claim without
converting those conditions into
independent features of the existent.

\subsection{The Conditional Character of Invariance}

The conditions established above
demonstrate that an invariance claim
is necessarily conditional.

This conditional character does not
mean that invariance is subjective
or that preservation depends upon
an arbitrary act of description.

It means that a claim of preservation
has determinate meaning only in
relation to the feature concerned,
the transformation considered,
the domain of applicability, and
the criterion by which preservation
is assessed.

The analytical specification of
these conditions does not create
the invariant. It establishes what
is being claimed and the circumstances
under which the claim may be justified.

The distinction between analytical
specification and ontological grounding
must therefore be maintained.

Analytical specification concerns
the formulation and evaluation of
the claim. Ontological grounding
concerns whether the feature identified
by that claim is genuinely preserved
in the relevant existent or relation.

Neither should be substituted for
the other.

An invariance claim may be precisely
formulated and nevertheless fail
because its proposed preservation
does not obtain. Conversely, a
feature may be preserved even when
an adequate account of that preservation
has not yet been formulated.

The distinction prevents the conditions
of description from being mistaken
for the conditions of existence.

\subsection{Synthesis}

An invariant cannot be adequately
defined as something that simply
remains unchanged.

A meaningful invariance claim requires
the identification of what is claimed
to be preserved, the transformation
under consideration, the domain
within which the claim applies,
the respect in which preservation
is assessed, and the criterion
by which preservation is established.

The admissibility of the transformation
and the scope of validity delimit
the applicability of the claim.
The distinction between formal
and ontological invariance determines
whether a formally established
preservation warrants an ontological
interpretation.

These conditions clarify the meaning
and limits of invariance without
presupposing that an invariant
necessarily exists.

They also prevent resemblance,
approximate correspondence, persistence,
continuity, and formal preservation
from being treated as interchangeable
with ontological invariance.

The framework developed here establishes
how claims of invariance are to be
specified and evaluated. It does
not establish that every existent
possesses an invariant, that every
transformation preserves something,
or that preservation is unlimited
across domains and respects.

The next section examines the relation
between invariance and transformation
more closely, particularly the possible
coexistence of preservation and
alteration within a single transformation.

% =============================================================================
% END OF SECTION 03 OF CHAPTER 02
% =============================================================================

% =============================================================================
% SECTION 04 — INVARIANCE AND TRANSFORMATION
% =============================================================================

% =============================================================================
% Project : THE_PHILOSOPHY_OF_VIABILITY
% Book    : Book I — Ontology of Viability
% Book ID : VS.B01
%
% File    :
%   2_MAINMATTER/PART_I_IDENTITY_AND_MANIFESTATION/
%   CHAPTER_02_INVARIANCE/sections/
%   section_04_invariance_and_transformation.tex
%
% Section : 04 — Invariance and Transformation
% Chapter : 02 — Invariance
%
% Purpose :
%   To examine the ontological and analytical relation between
%   transformation, preservation, and alteration, establishing
%   the conditions under which invariance may be attributed to
%   a specified aspect of an existent undergoing transformation.
%
% Conceptual Scope :
%   Transformation, preservation, alteration, partial preservation,
%   total preservation, configurations of determinations,
%   relational organization, and the limits of invariance claims.
%
% =============================================================================

\section{Invariance and Transformation}
\label{sec:invariance-and-transformation}

\subsection{The Relation Between Invariance and Transformation}

Invariance and transformation are sometimes treated as
opposing conditions. According to this interpretation,
whatever remains invariant must remain unchanged, whereas
whatever undergoes transformation must necessarily lose
its previous determinations.

This opposition is inadequate because it treats change
and preservation as properties that must characterize
an existent uniformly and in its entirety.

A transformation need not affect every determination,
relation, or organizational feature of an existent in
the same manner. Some aspects may be altered, while others
may remain preserved relative to the transformation
under consideration.

Conversely, the preservation of a particular determination
or relation does not establish that the existent has
remained unchanged in every other respect.

Transformation and invariance therefore belong to
analytically distinguishable but potentially compatible
descriptions of an existent and its determinations.

Their relation cannot be adequately examined through
the undifferentiated opposition between change and
unchange. It requires identifying what is transformed,
what is assessed as preserved, and the respect in which
preservation or alteration is attributed.

The central question is consequently not whether an
existent is either transformed or invariant without
qualification. It is whether a specified feature,
determination, or relation remains preserved under
a specified transformation, within a specified domain,
and according to an appropriate criterion.

Invariance does not negate transformation, and
transformation does not, merely by occurring, negate
every possibility of invariance.

\subsection{Transformation Does Not Necessarily Eliminate Invariance}

A transformation involves an alteration in at least
one relevant aspect of a configuration, determination,
or relation, such that the resulting condition differs
from the prior condition in the respect under examination.

However, the occurrence of such an alteration does not
establish that every other aspect has also been altered.

A transformation may affect certain determinations
while leaving others preserved. It may modify quantitative
properties without eliminating a specified structural
relation. It may alter the properties of relata while
leaving a particular relational organization intact.

In these cases, the transformation and the preservation
of a specified feature are compatible.

The preservation claim must nevertheless identify
the feature whose invariance is being assessed. It is
not sufficient to observe that something appears to
have remained unchanged. The relevant determination
or relation must be specified, and its preservation
must be evaluated under the transformation in question.

The distinction is important because the existence
of alteration in one respect provides no sufficient
grounds for denying preservation in another.

A transformation that alters the quantitative properties
of certain constituents, for example, does not thereby
establish that every relation among those constituents
has also been altered.

Likewise, a transformation that changes the material
constitution of an existent does not, without further
analysis, determine whether its structural organization
has been preserved or altered.

The possibility of invariance within transformation
therefore depends on the differentiation of the aspects
under examination. It does not require that the existent
remain wholly unaffected.

\subsection{Invariance Does Not Imply the Absence of Transformation}

The converse limitation is equally necessary.

The preservation of a determination or relation does
not establish that no transformation has occurred.

An invariant may be identified within a configuration
that has undergone alteration in other respects.
The preservation of one feature cannot be generalized
to the entire existent unless the scope of the claim
explicitly warrants such a generalization.

For example, the preservation of a specified relational
pattern does not establish that every participating
determination has remained unchanged. Similarly, the
preservation of a structural feature does not establish
that the structure has remained identical in every
respect or that no other aspect of its organization
has been modified.

The error arises when a limited preservation claim
is extended beyond its subject and criterion.

An invariance claim concerns the preservation of
a specified feature under specified conditions. It
does not automatically characterize the total condition
of the existent.

Accordingly, the assertion that a particular determination
remains invariant must not be converted into the assertion
that the existent as a whole has undergone no transformation.

The presence of an invariant therefore does not establish
immutability. It establishes, at most, the preservation
of the specified feature within the scope of the
transformation and criterion under consideration.

\subsection{The Coexistence of Preservation and Alteration}

Preservation and alteration may coexist within the
same transformation.

This possibility arises because the determinations
and relations associated with an existent need not
be affected uniformly.

A transformation may preserve one determination,
alter another, modify a relation, and establish a
different configuration among other determinations.

These outcomes are not mutually exclusive merely
because they occur within the same transformed
configuration.

Their compatibility depends on the respects in
which preservation and alteration are attributed.

A contradiction would arise if the same subject,
the same determination or relation, and the same
respect were asserted to be both preserved and altered
under the same transformation, within the same domain,
and according to the same criterion.

Where the relevant subjects or respects differ,
the two attributions may consistently coexist.

For example, a particular structural relation may
remain preserved while the quantitative properties
of the elements participating in that relation are
altered. The relational feature and the quantitative
properties are not the same subject of assessment.
Their different outcomes therefore do not constitute
a contradiction.

Likewise, the preservation of an organizational
pattern may coexist with the alteration of the
determinations occupying particular positions within
that pattern.

The distinction does not mean that every apparent
coexistence of preservation and alteration is coherent.
The subject, transformation, domain, respect, and
criterion must be sufficiently specified to determine
whether the two claims concern compatible aspects
or make incompatible assertions about the same one.

The coexistence of preservation and alteration is
therefore a possibility that requires analytical
specification, not a conclusion that can be assumed
in every case.

\subsection{Partial and Total Preservation}

The distinction between partial and total preservation
concerns the extent of what is preserved relative to
a specified assessment.

Partial preservation occurs when some of the features
included in the assessment remain preserved while
others are altered.

It does not necessarily indicate that the transformation
is minor, nor does it determine the ontological
significance of the features that remain preserved.
A transformation may be extensive in several respects
while preserving a particular determination or relation.

Total preservation requires a more restricted and
carefully specified meaning.

It can be attributed only when every feature included
within the explicitly defined scope of assessment
remains preserved according to the relevant criterion.

The qualification concerning scope is indispensable.
No finite or restricted assessment can establish
that every aspect of an existent, without qualification,
has remained unchanged.

Consequently, total preservation within a specified
assessment must not be confused with absolute
immutability or the absence of transformation in
every conceivable respect.

Even when all features included in an assessment
remain preserved, the assessment may exclude other
determinations, relations, or aspects of the existent.
The conclusion must therefore remain proportionate
to the scope of the examination.

Partial and total preservation are thus comparative
descriptions of the extent to which specified features
are preserved. Neither description, by itself,
establishes the numerical identity of the existent.

Nor does the distinction determine whether the
preserved features are ontologically fundamental,
structurally significant, or merely restricted
characteristics. Such judgments require additional
arguments and cannot be derived from the extent
of preservation alone.

\subsection{Transformation and the Configuration of Determinations}

Transformation may affect individual determinations
as well as the configuration in which those
determinations are organized.

These two levels of assessment must remain distinct.

A determination may remain preserved while its
position or role within a configuration changes.
Conversely, a configuration may retain a specified
organizational relation even when some of the
determinations participating in it undergo alteration.

The preservation of an individual determination
therefore does not automatically establish the
preservation of the configuration in which it occurs.

Similarly, the preservation of an organizational
relation does not require that every determination
participating in that relation remain unchanged.

The distinction becomes especially important when
the configuration itself is the subject of an
invariance claim.

In such cases, the analysis must identify whether
the relevant feature concerns:

\begin{enumerate}
    \item the preservation of particular determinations;
    \item the preservation of relations among determinations;
    \item the preservation of a specified organization
    of determinations and relations; or
    \item the preservation of a combination of these
    features under the transformation in question.
\end{enumerate}

These are distinct possible subjects of assessment.
They must not be treated as interchangeable.

A configuration may undergo alteration even when
some of its constituent determinations remain
preserved. It may also retain a specified organizational
feature despite changes in some of its constituents.

Furthermore, the preservation of a relation does
not necessarily entail the preservation of the
determinations participating in it. Whether a
particular relation can remain preserved under
such changes depends on the nature of the relation
and the criterion by which its preservation is assessed.

The analysis of invariance must therefore distinguish
between the preservation of determinations, the
preservation of relations, and the preservation
of configurations.

No one of these forms of preservation should be
assumed to establish the others without further
justification.

\subsection{Transformation and the Limits of Generalization}

An invariant established under a particular
transformation cannot automatically be generalized
to other transformations.

A determination or relation may remain preserved
under one transformation while undergoing alteration
under another. The same feature may also be preserved
within a restricted domain but fail to remain
preserved when the transformation extends beyond
that domain.

The validity of an invariance claim must therefore
remain bounded by the conditions under which it
has been established.

These conditions include the specified subject,
the transformation or admissible class of transformations,
the domain of applicability, the respect of preservation,
and the criterion used to assess the result.

This limitation does not make invariance arbitrary.
It establishes the conditions under which a claim
is warranted and prevents a restricted instance
of preservation from being treated as an unrestricted
ontological principle.

A feature demonstrated to remain preserved under
a defined class of transformations may legitimately
be described as invariant relative to that class,
provided that the relevant conditions and limits
have been established.

Such a conclusion does not establish that the
feature is preserved under every possible
transformation, nor does it establish that the
same result holds in every domain or respect.

The distinction between restricted and unrestricted
claims is therefore essential.

An unrestricted claim would require a correspondingly
unrestricted justification. It cannot be inferred
merely from the successful establishment of invariance
in a limited set of cases.

The scope of an invariance claim must consequently
remain proportionate to the scope of the argument
supporting it.

\subsection{Transformation, Invariance, and the Question of Identity}

The coexistence of preservation and alteration
raises a further question concerning the relation
between transformation and identity.

The preservation of a determination or relation
across a transformation may contribute to the
characterization of the existent and the manner
in which its configuration has changed.

Nevertheless, the existence of an invariant does
not, by itself, establish that the existent considered
before and after the transformation is numerically
identical.

An invariant specifies a feature that remains
preserved under defined conditions. Numerical
identity concerns a different ontological question.
The former cannot be substituted for an argument
establishing the latter.

The converse limitation is equally important.

The alteration or disappearance of a previously
identified invariant does not, by itself, establish
the disappearance of the identity of the existent.

The loss of a particular determination, relation,
or organizational feature may be relevant to an
investigation of identity, but its ontological
significance cannot be determined solely by the
fact that the feature has ceased to be preserved.

These limitations follow from the analytical
distinction between identity and invariance.
Neither concept can be reduced to the other, and
neither can independently settle every question
concerning the other.

Transformation may therefore be examined in terms
of preserved and altered features without presuming
that these findings resolve the identity of the
existent.

The question of what remains invariant and the
question of whether the existent remains numerically
identical must remain distinct, even when their
investigation is related.

\subsection{Analytical Possibilities of Preservation and Alteration}

The relation between transformation, preservation,
and alteration can be clarified through several
analytical possibilities.

A transformation may involve the preservation of
a specified determination or relation while other
aspects remain outside the scope of assessment.

It may involve the alteration of a specified
determination or relation without any invariant
having been established.

It may also involve the coexistence of preservation
and alteration, provided that the respective claims
concern compatible subjects or respects.

These possibilities clarify the logical relations
between transformation and invariance. They do
not constitute an exhaustive classification of
transformations, nor do they imply that every
transformation must produce an identifiable invariant.

The analytical distinctions can be stated as follows:

\begin{enumerate}
    \item \textbf{Preservation under transformation:}
    A specified feature remains preserved according
    to a defined criterion, notwithstanding alteration
    elsewhere in the configuration.

    \item \textbf{Alteration under transformation:}
    A specified feature undergoes alteration. This
    finding does not, by itself, determine whether
    other features remain preserved.

    \item \textbf{Coexisting preservation and alteration:}
    Some specified features remain preserved while
    others undergo alteration within the same
    transformation.

    \item \textbf{Undetermined preservation:}
    A transformation is established, but the available
    analysis does not determine whether a particular
    feature has been preserved, altered, or replaced.
\end{enumerate}

The fourth possibility is important because the
absence of an established invariant must not be
confused with the establishment of alteration.
An insufficiently determined case remains
unresolved with respect to the relevant preservation
claim.

Likewise, the identification of alteration in one
respect does not establish alteration in every
other respect.

These possibilities therefore distinguish what
may be asserted from what remains undetermined.
They do not establish a universal sequence or
a necessary outcome for every transformation.

The analytical value of these distinctions lies
in preventing the reduction of transformation
to either complete preservation or complete
alteration.

\subsection{Synthesis}

Invariance and transformation are not mutually
exclusive conditions.

A transformation may preserve certain determinations,
relations, or organizational features while altering
others. The occurrence of alteration in one respect
does not necessarily eliminate invariance in another,
and the preservation of a particular feature does
not establish the absence of transformation in
the existent as a whole.

The distinction between partial and total preservation
clarifies the extent of a preservation claim.
Total preservation can be asserted only relative
to an explicitly specified scope of assessment;
it must not be confused with absolute immutability.

The analysis must also distinguish between the
preservation of individual determinations, relations
among determinations, and the organization of a
configuration. These forms of preservation are
related but not interchangeable.

Every invariance claim remains bounded by its
specified subject, transformation, domain, respect,
and criterion. Its conclusions cannot be generalized
beyond the conditions under which preservation
has been established.

Finally, neither the preservation nor the alteration
of particular features is sufficient, by itself,
to establish or negate the numerical identity
of an existent.

Transformation and invariance therefore provide
distinct but related perspectives for examining
what is preserved and what is altered. Their
analysis contributes to the characterization of
an existent without replacing the separate
ontological question of its identity.

The next section examines the relation between
invariance and difference, addressing how the
preservation of specified features may coexist
with difference within and across transformed
configurations.

% =============================================================================
% END OF SECTION 04 OF CHAPTER 02
% =============================================================================

% =============================================================================
% SECTION 05 — INVARIANCE AND DIFFERENCE
% =============================================================================


% =============================================================================
% Section 05 — Invariance and Difference
% Book I — Ontology of Viability
% Chapter 02 — Invariance
%
% Purpose:
%   To examine the relation between invariance and difference,
%   establishing how preservation may coexist with difference
%   within and across transformed configurations.
%
% Conceptual Policy:
%   1. Invariance does not eliminate difference.
%   2. Difference does not necessarily eliminate invariance.
%   3. Preservation and difference must be assessed with respect
%      to specified determinations, relations, and conditions.
%   4. Similarity, equivalence, and identity must not be conflated.
%   5. The persistence of an invariant does not establish
%      the absence of ontological difference.
%   6. Difference must not be reduced to variation in a single
%      identifiable feature.
%
% =============================================================================

\section{Invariance and Difference}
\label{sec:invariance-and-difference}

\subsection{The Ontological Problem of Invariance and Difference}

The relation between invariance and difference constitutes
a further problem in the ontological analysis of preservation.

Invariance concerns the preservation of a specified
determination or relation across a transformation.
Difference concerns the distinction or non-identity
between determinations, configurations, or existents.

These concepts address different aspects of ontological
reality and must not be treated as mutually exclusive.

The preservation of a particular feature does not
necessarily eliminate differences between the configurations
in which that feature obtains. Similarly, the existence
of difference does not necessarily imply that every
determination or relation has been altered.

An existent may exhibit invariant features while
undergoing differences in other respects. Different
existents may also share particular invariant features
without thereby becoming numerically identical.

The central problem is therefore to clarify how preservation
and difference can coexist without reducing either
concept to the other.

\subsection{Invariance Does Not Eliminate Difference}

The preservation of a determination or relation does
not establish the absence of difference.

A particular feature may remain preserved while other
determinations, relations, or configurations undergo
alteration.

Consequently, two configurations may share an invariant
and nevertheless differ in several other respects.

The preservation of a structural organization, for example,
does not necessarily imply the preservation of every
quantitative or qualitative determination associated
with that organization.

Likewise, the preservation of a particular relation
does not establish that the relata participating
in that relation are identical in every respect.

Invariance must therefore be understood as the preservation
of a specified feature rather than as the elimination
of all difference.

To infer the absence of difference from the existence
of an invariant would be to extend a restricted
preservation claim beyond its legitimate scope.

\subsection{Difference Does Not Necessarily Eliminate Invariance}

The converse is equally important.

The existence of difference between two configurations
does not necessarily establish that no invariant
is preserved between them.

A transformation may produce differences in certain
determinations while leaving other features intact.

The configurations before and after the transformation
may therefore be different without being entirely
discontinuous in every relevant respect.

The same possibility applies to relations between
different existents.

Two existents may differ in their determinations,
compositions, or configurations while sharing a
particular relational or structural feature.

Such preservation does not eliminate their differences.
It establishes only that a specified feature remains
common or preserved under the conditions considered.

Difference and invariance must consequently be examined
at the appropriate analytical level.

\subsection{Difference Between Determinations and Difference Between Configurations}

The concept of difference must be considered at more
than one analytical level.

A difference between individual determinations is not
necessarily equivalent to a difference between entire
configurations.

Two configurations may contain different determinations
while preserving a particular organizational relation.

Conversely, two configurations may contain similar
determinations while differing in the way those
determinations are organized or related.

The preservation of individual determinations therefore
does not automatically establish the preservation
of a configuration.

Likewise, the alteration of individual determinations
does not necessarily establish the alteration
of every structural or relational feature.

An adequate account of invariance must distinguish
between the preservation or alteration of individual
determinations and the preservation or alteration
of their configurations.

Without this distinction, different forms of preservation
and difference may be incorrectly conflated.

\subsection{Difference and the Scope of Comparison}

Every assertion of difference presupposes a basis
of comparison.

The relevant basis may concern determinations,
relations, structures, compositions, or other
specified aspects of existents.

Two configurations may be different under one criterion
while remaining equivalent under another.

For example, two structures may differ in their
constituent elements while preserving an equivalent
organizational pattern.

Such equivalence does not establish numerical identity.
It establishes a correspondence under a specified
criterion of comparison.

Similarly, the identification of a difference in
one respect does not establish that the compared
configurations differ in every respect.

The scope of comparison must therefore be specified
before conclusions about difference or preservation
can be established.

This requirement parallels the conditions developed
for identifying invariants, since both concepts
require a determinate subject and respect of analysis.

\subsection{Invariance, Similarity, and Equivalence}

Invariance must also be distinguished from similarity
and equivalence.

Similarity indicates a degree of resemblance between
determinations or configurations.

Equivalence indicates correspondence under a specified
criterion or relation of comparison.

Invariance, by contrast, concerns the preservation
of a specified feature across a transformation
or within a defined domain of comparison.

These concepts may be related, but they are not
interchangeable.

Two configurations may be similar without preserving
the same determination through a transformation.
Two structures may be equivalent under a formal
criterion without being ontologically identical.

Likewise, the preservation of an invariant does
not necessarily imply that the initial and resulting
configurations are similar in every respect.

The distinction is particularly important when
formal relations are used to describe ontological
configurations.

A formal equivalence may establish a correspondence
within a model, but it does not automatically
establish the numerical identity of the existents
represented by that model.

The ontological status of the compared configurations
must therefore remain distinct from the formal
criterion through which their correspondence
is established.

\subsection{Invariance Across Different Existents}

The problem of invariance and difference is not
restricted to transformations occurring within
a single existent.

A particular determination or relation may be
present in different existents.

The presence of a shared feature does not establish
that the existents possessing it are numerically
identical.

Rather, it establishes that a particular feature
or relation is common to them under the relevant
criterion.

The distinction between shared invariance and
numerical identity is essential.

Different existents may exhibit similar structures,
preserve equivalent relations, or possess common
determinations while remaining distinct existents.

The existence of common features therefore cannot
be used as a sufficient basis for reducing
ontological difference to identity.

This also indicates that invariance may be relevant
to comparisons among existents without determining
their numerical identity.

\subsection{Difference Without the Loss of Every Invariant}

The occurrence of difference does not necessarily
eliminate every invariant associated with an existent.

A transformation may alter a particular determination
while preserving a relation or structural feature.

Similarly, the emergence of a new configuration
may involve both the alteration of certain features
and the preservation of others.

The identification of difference must therefore
not be treated as sufficient grounds for denying
all forms of invariance.

However, the converse limitation must also be maintained:
the preservation of one feature cannot establish
that every other feature remains unchanged.

The coexistence of difference and invariance requires
an analysis that remains specific to the relevant
determinations, relations, and conditions.

Neither complete preservation nor complete alteration
should be presumed without sufficient ontological
justification.

\subsection{The Limits of Invariance as a Basis of Distinction}

Although invariant features may contribute to
the differentiation or comparison of existents,
they cannot be treated as an exhaustive basis
for establishing ontological distinction.

Two existents may share particular invariant
features while differing in other determinations
or relations.

Conversely, an existent may undergo alteration
in previously identified invariant features
without thereby losing its identity.

An invariant can therefore provide a basis for
describing a particular aspect of an existent,
but it cannot independently determine the full
ontological status of that existent.

This limitation is consistent with the distinction
between identity and invariance established earlier.

The analysis of difference must consequently
remain open to determinations and relations
that are not captured by a particular invariance
claim.

\subsection{Synthesis}

Invariance and difference are not mutually exclusive
ontological conditions.

The preservation of a specified determination
or relation does not eliminate differences in
other respects. Similarly, the existence of
difference does not necessarily eliminate every
invariant.

Difference may concern individual determinations,
relations, or entire configurations. Its assessment
therefore requires a specified basis and scope
of comparison.

Similarity and equivalence must also remain
distinguishable from invariance and numerical
identity. A shared or preserved feature does
not establish that the existents possessing
it are identical.

Invariance contributes to the analysis of
preservation and comparison, but it cannot
exhaust the ontological problem of difference.

The next section examines the relation between
invariance and structure, addressing whether
the preservation of structural organization
can be distinguished from the preservation
of individual determinations and relations.


% =============================================================================
% SECTION 06 — INVARIANCE AND STRUCTURE
% =============================================================================

% =============================================================================
% Project : THE_PHILOSOPHY_OF_VIABILITY
% Book    : Book I — Ontology of Viability
% Book ID : VS.B01
%
% File    :
%   2_MAINMATTER/PART_I_IDENTITY_AND_MANIFESTATION/
%   CHAPTER_02_INVARIANCE/sections/
%   section_06_invariance_and_structure.tex
%
% Section : 06 — Invariance and Structure
% Chapter : 02 — Invariance
%
% Purpose :
%   To examine structural invariance as a specified form of
%   preservation, distinguishing the preservation of components,
%   relations, and organizational arrangements from the numerical
%   identity of the existent in which they obtain.
%
% Conceptual Scope :
%   Structure, structural determination, component variation,
%   relational organization, structural transformation,
%   structural equivalence, and the limits of structural explanation.
%
% =============================================================================

\section{Invariance and Structure}
\label{sec:invariance-and-structure}

\subsection{The Ontological Relation Between Structure and Invariance}

Structure and invariance are related because an existent
may retain a specified structural determination across
a transformation. Changes in components, relations,
configurations, or material constitution do not necessarily
alter every aspect of the organization under consideration.

The preservation of a specified structural determination
under a defined transformation constitutes a particular
form of invariance, which may be described as structural
invariance.

Structural invariance concerns the preservation of an
identified organizational feature under specified
conditions. It does not refer simply to the continued
existence of the subject, nor does it establish that
the existent remains numerically identical through
every transformation under consideration.

Structure and identity address different ontological
questions. Structure concerns an organization or
relational determination of an existent. Identity
concerns the existent as such and cannot be reduced
to the preservation of any particular organizational
arrangement.

A structural determination may remain preserved
across a transformation without establishing the
numerical identity of the subjects before and after
that transformation. Conversely, an existent may
undergo structural alteration without that alteration
alone establishing the loss of its identity.

The distinction is therefore not between structure
and an existent considered as though they were
independent objects. It is between a particular
determination of an existent and the ontological
question of the identity of that existent.

Structural invariance must consequently be examined
as a determinate form of invariance, not as a general
explanation or criterion of identity.

\subsection{Structure as an Organized Determination}

A structure is not adequately characterized as
a mere aggregate of components.

An aggregate may be considered in terms of the
elements it contains, whereas structural determination
concerns the organization and relations through
which components or other determinations constitute
an arrangement.

The structural character of an existent may involve:

\begin{itemize}
    \item the organization of its components;
    \item the relations connecting those components;
    \item the arrangement of relations within a
          specified configuration;
    \item the constraints governing admissible
          configurations;
    \item the dependencies through which particular
          organizational determinations are maintained.
\end{itemize}

These aspects are analytically distinguishable.
They need not remain equally preserved under
a transformation.

A transformation may alter particular components
while preserving some relations among them. It
may preserve a relational arrangement while
altering other organizational constraints.
It may also modify the configuration as a whole
while retaining a more restricted structural feature.

For this reason, the assertion that a structure
remains invariant requires specification.

It is insufficient to state that ``the structure
remains the same'' without identifying the
organizational determination being considered,
the transformation under examination, and the
criterion according to which preservation is
established.

Structural invariance is therefore relative
to a specified structural respect. It cannot
be attributed to an entire structure without
clarifying which of its aspects are included
in the claim.

\subsection{Structural Preservation and Component Variation}

One possible form of structural invariance occurs
when components change while a specified
organizational relation remains preserved.

The elements occupying particular structural
positions may be modified or replaced without
necessarily altering the organization assessed
under a defined criterion.

For example, two relational configurations may
contain different elements while preserving
a corresponding arrangement of relations.
In such a case, the preservation concerns
the specified relational organization, not
the numerical identity of the components.

This possibility demonstrates that structural
invariance cannot be reduced to the persistence
of individual components.

However, the converse must also be recognized.

The persistence of components does not guarantee
the preservation of their structural organization.
The same components may enter into different
relations, occupy different positions, or become
organized according to a different arrangement.

Component preservation and structural preservation
are therefore analytically distinct.

Neither can be inferred automatically from the
other.

Furthermore, the replacement of a component
does not necessarily preserve the relevant
structure. Whether the organization remains
preserved depends on the structural respect
and criterion under consideration.

The fact that different components occupy
corresponding positions is insufficient unless
the relevant relations and organizational
conditions are also specified.

Structural invariance must consequently be
established through the assessment of the
organization itself, rather than inferred
solely from the persistence or replacement
of its components.

\subsection{Structural Relations and Organizational Preservation}

The preservation of a structure may concern
more than the persistence of individual
relations considered separately.

An organization may depend on how several
relations are connected, ordered, or mutually
constrained.

Consequently, the preservation of individual
relations does not necessarily establish
the preservation of their overall organization.

A configuration may retain several relations
while altering the arrangement through which
those relations constitute a structure.

Conversely, a specified organizational
arrangement may remain preserved even when
some individual relations are realized through
different components or determinations.

The distinction requires separating at least
three possible subjects of structural assessment:

\begin{enumerate}
    \item \textbf{Individual relations:}
    Whether particular relations remain preserved
    under the transformation.

    \item \textbf{Relational organization:}
    Whether the arrangement and interdependence
    of relations remain preserved according
    to a specified criterion.

    \item \textbf{Structural constraints:}
    Whether the conditions restricting or
    governing admissible configurations remain
    preserved.
\end{enumerate}

These subjects may be related, but they are
not interchangeable.

The preservation of individual relations
does not automatically establish the
preservation of their organization. Likewise,
the preservation of a structural constraint
does not necessarily establish that every
relation governed by that constraint remains
unchanged.

A structural invariance claim must therefore
identify whether it concerns individual
relations, their organization, particular
constraints, or a specified combination
of these features.

Without this specification, the assertion
of structural preservation remains
insufficiently determined.

\subsection{Structural Transformation and the Limits of Preservation}

Structural transformations may occur at
different analytical levels and with
different consequences for the organization
under consideration.

A transformation may:

\begin{enumerate}
    \item alter individual components while
          preserving a specified organizational
          relation;

    \item alter some relations while preserving
          a broader structural arrangement;

    \item reorganize relations while retaining
          particular structural constraints;

    \item alter the organizational arrangement
          to such an extent that the previously
          specified structural feature is no
          longer preserved.
\end{enumerate}

These possibilities indicate that structural
change is not a single, uniform phenomenon.

A structure may remain invariant with respect
to one transformation and fail to remain
invariant with respect to another. It may
also preserve one structural respect while
undergoing alteration in another.

Thus, structural invariance cannot be asserted
without specifying the transformation or
class of transformations, the relevant domain,
the structural respect, and the criterion
of preservation.

The recognition of structural change also
does not establish, by itself, that the
identity of the existent has been lost.

Such a conclusion would require an independent
ontological argument concerning identity.
The alteration of an organizational arrangement
may be relevant to that investigation, but
it cannot independently settle it.

Likewise, the preservation of a structural
arrangement cannot establish that the same
existent has persisted. A corresponding
or equivalent organization may obtain in
distinct existents.

Structural transformation and structural
preservation must therefore be examined
without presuming that either determines
the numerical identity of the existent.

\subsection{Structural Equivalence and Numerical Identity}

A particularly important distinction concerns
structural equivalence and numerical identity.

Two existents may exhibit corresponding
structural organizations without being
numerically identical.

Their structures may be equivalent under
a specified criterion, while the existents
in which those structures obtain remain
distinct.

Structural equivalence expresses a relation
between structural determinations. Numerical
identity concerns whether an existent is
one and the same existent.

These are different ontological questions.

The distinction prevents the reduction
of identity to structure.

If structural equivalence were treated
as sufficient for numerical identity,
distinct existents possessing corresponding
organizations would incorrectly be treated
as one existent.

The same problem arises when the preservation
of a structure through transformation is
taken as proof of the persistence of the
existent associated with it.

Such an inference exceeds what structural
invariance alone can establish.

Structural equivalence may support a
comparison between configurations, and
structural invariance may establish the
preservation of a specified organization.
Neither conclusion independently settles
the numerical identity of the existents
being compared.

The ontological status of the subjects
must therefore remain distinct from the
equivalence or invariance of their
structural determinations.

\subsection{Structural Sameness and the Limits of Structural Identity}

Expressions such as ``the same structure''
or ``structural identity'' require
qualification.

They may refer to the preservation of
an organizational arrangement across
a transformation, or to the correspondence
of two structures under a specified
criterion.

In either case, the expression concerns
the sameness or equivalence of structural
determinations. It must not be confused
with the ontological identity of the
existent whose structure is being examined.

A structure may be treated as the same
under a defined criterion even when
the components occupying its positions
have changed.

Likewise, two structures may be described
as equivalent under a formal criterion
while belonging to distinct existents.

Such descriptions establish a particular
analytical relation. They do not establish
that the associated existents are
numerically one and the same.

Moreover, an existent must not be assumed
to be exhausted by a single structural
determination or by one structural
description.

Different structural respects may be
relevant to the same existent. A
transformation may preserve one of
these respects while altering another.

Consequently, any assertion of structural
sameness must specify its subject,
criterion, and scope.

The expression cannot serve as an
unqualified substitute for ontological
identity, nor can a particular structural
description be presumed to exhaust
the determinations of the existent.

\subsection{The Limits of Structural Explanation of Identity}

The preceding analysis establishes
a necessary limitation: structural
invariance is a possible form of
preservation, but the identity of
an existent cannot be reduced to
the preservation of its structure.

This limitation operates in both
directions.

First, structural preservation does
not prove numerical identity.

Equivalent structures may belong
to distinct existents, and the
preservation of an organizational
arrangement does not, by itself,
establish the persistence of the
existent associated with it.

Second, structural transformation
does not prove the loss of identity.

An existent may undergo changes
in its organization without the
mere occurrence of those changes
determining that it has become
another existent.

Third, the absence of an established
structural invariant does not establish
the absence of identity.

The failure to identify a preserved
structural feature may indicate
that the relevant structure has
changed, that the criterion is
inadequate, or that the available
analysis does not determine the
matter. It cannot independently
settle the ontological question
of identity.

Structure may therefore contribute
to the analysis of an existent
without becoming an exhaustive
ground of its identity.

Identity and structure must remain
analytically distinct. The fact
that structure is a determination
of an existent does not make it
identical with the existent's
identity, nor does it establish
that structural preservation is
a necessary or sufficient condition
of numerical identity.

The proper conclusion is not that
structure is ontologically irrelevant
to identity, but that its relevance
must not be converted into reduction.

\subsection{Synthesis}

Structural invariance concerns the
preservation of a specified organizational
determination across a defined domain
of transformation.

Its analysis requires distinguishing:

\begin{itemize}
    \item the preservation of components
          from the preservation of their
          organization;

    \item the preservation of individual
          relations from the preservation
          of their relational arrangement;

    \item the preservation of structural
          constraints from the preservation
          of every configuration governed
          by those constraints;

    \item structural equivalence from
          numerical identity;

    \item structural transformation from
          the loss of ontological identity;

    \item structural sameness under a
          specified criterion from the
          identity of the existent itself.
\end{itemize}

The preservation of structure is neither
a sufficient condition for numerical
identity nor a universally necessary
condition that can be imposed upon
every existent.

Structural invariance remains an
important analytical category within
the ontology of invariance. Its validity,
however, is restricted to the structural
determination, transformation, domain,
and criterion of preservation under
consideration.

The identity of an existent must
consequently remain an independent
ontological question, irreducible
to the invariance, equivalence, or
transformation of its structure.

The next section examines the relation
between invariance and composition,
addressing the distinction between
the preservation of an organized
arrangement and the preservation
or alteration of the constituents
through which that arrangement
is constituted.

% =============================================================================
% END OF SECTION 06 OF CHAPTER 02
% =============================================================================

% =============================================================================
% SECTION 07 — INVARIANCE AND COMPOSITION
% =============================================================================

% =============================================================================
% Section 07 — Invariance and Composition
% Book I — Ontology of Viability
% Chapter 02 — Invariance
%
% Purpose:
%   To examine the relation between invariance and composition,
%   distinguishing the preservation of compositional determinations
%   from structural invariance and ontological identity.
%
% Conceptual Policy:
%   1. Compositional invariance is not identical to structural
%      invariance.
%   2. The persistence of components does not necessarily establish
%      the persistence of their composition.
%   3. Compositional change does not by itself establish the loss
%      of ontological identity.
%   4. The preservation of composition does not establish numerical
%      identity.
% =============================================================================

\section{Invariance and Composition}
\label{sec:invariance-and-composition}

\subsection{The Ontological Problem of Composition}

Composition concerns the manner in which multiple components
constitute or participate in the constitution of an existent.
It addresses not merely the presence of components, but also
their compositional arrangement and the determinate relations
through which they belong to a composed whole.

The relation between composition and invariance becomes
significant when an existent undergoes transformation while
some aspect of its composition remains preserved.

A composed existent may retain the same components while
their arrangement changes. It may also undergo the replacement
of some components while preserving a specified compositional
organization. In other cases, both components and their
compositional relations may change.

These possibilities indicate that compositional preservation
cannot be identified simply with the persistence of components
or with the persistence of the whole.

Compositional invariance must therefore be examined as a
specific form of invariance concerning the constitution of
an existent.

It does not, by itself, provide an exhaustive account of
the identity of that existent.

\subsection{Composition and the Presence of Components}

The presence of components is one aspect of composition,
but it is not sufficient to characterize composition as such.

A collection of components does not necessarily constitute
the same composed arrangement under every possible relation
or organization.

The same components may participate in different compositions.
Conversely, different components may occupy comparable
compositional positions within an organization.

This distinction requires separating two questions:

\begin{enumerate}
    \item Which components are present in a composed existent?
    \item How are those components compositionally constituted
          as part of that existent?
\end{enumerate}

The first question concerns component membership.
The second concerns compositional organization.

A claim of compositional invariance may concern either
of these respects or both. Consequently, the precise respect
of preservation must be specified.

The persistence of component membership cannot automatically
be treated as the persistence of composition. Likewise,
the replacement of a component does not necessarily imply
that every compositional determination has been altered.

\subsection{Compositional Invariance}

Compositional invariance refers to the preservation of a
specified compositional determination across a defined
transformation.

Such preservation may concern:

\begin{itemize}
    \item the membership of components within a composed whole;
    \item the number of components;
    \item the distribution of components among compositional
          positions;
    \item the relations through which components participate
          in a composition;
    \item the organization of components into a determinate
          compositional arrangement.
\end{itemize}

These respects are analytically distinguishable.
They need not be preserved together, and the preservation
of one does not establish the preservation of the others.

For example, the number of components may remain constant
while their membership changes. Component membership may
remain unchanged while their arrangement is reorganized.
A compositional arrangement may also remain equivalent
under a specified criterion even when particular components
are replaced.

Therefore, the assertion that composition remains invariant
requires a clear account of what aspect of composition
is being preserved.

Without such specification, compositional invariance
remains indeterminate.

\subsection{Component Replacement and Compositional Preservation}

An important case occurs when one or more components are
replaced while the broader compositional organization
remains preserved.

In this situation, the identity of individual components
does not necessarily determine the invariance of the
composition as a whole.

A replacement may preserve the compositional role or
position previously occupied by a component. It may also
preserve the number and organization of components while
altering their particular constitution.

However, the preservation of compositional organization
must not be assumed merely because replacement has occurred
without an immediately visible alteration of the whole.

Whether composition remains invariant depends on the
criterion by which the composition is being evaluated.

If component identity is constitutive of the specified
compositional determination, replacement may constitute
a loss of that invariance. If the relevant respect concerns
only a broader organization, replacement may be compatible
with its preservation.

Thus, component replacement does not have one universally
determined consequence for compositional invariance.

Its significance depends on the compositional determination
under consideration.

\subsection{Compositional Reorganization}

Composition may also change without the replacement
or disappearance of its components.

The same components may be reorganized into a different
compositional arrangement. Their distribution, participation,
or relations may change while their individual presence
remains preserved.

This possibility demonstrates that component persistence
and compositional invariance are not equivalent.

Compositional reorganization may occur at different levels.
A limited alteration may affect only one part of a composition,
while a more extensive transformation may alter the
organization of the composed whole.

The degree and significance of such transformation cannot
be determined solely by counting the components that
remain unchanged.

A compositional determination may be preserved despite
some local modifications, or it may cease to be preserved
even when all individual components remain present.

The assessment therefore requires an explicit criterion
of compositional preservation.

\subsection{Composition and Structural Invariance}

Compositional invariance and structural invariance are
closely related, but they are not identical.

Compositional invariance concerns the preservation of
a specified determination relating to the constitution
of a composed whole through its components.

Structural invariance concerns the preservation of a
specified organizational pattern or relational arrangement.

The two may coincide in particular cases. A compositional
arrangement may be preserved through the preservation
of its structural organization. Nevertheless, their
analytical scopes remain distinguishable.

A composition may preserve its component membership
while undergoing structural reorganization. Conversely,
a structural organization may remain equivalent even
when some of its components are replaced.

This distinction prevents compositional invariance
from being reduced to structural invariance.

It also prevents the assumption that every structural
transformation necessarily constitutes a transformation
of composition in every relevant respect.

The relationship between the two must be determined
according to the particular compositional and structural
determinations involved.

\subsection{Compositional Change and the Identity of the Existent}

The transformation of composition raises a further
ontological question: whether a change in the composition
of an existent entails a change in its identity.

No universal answer can be derived from compositional
invariance alone.

The replacement of a component, the reorganization
of components, or the alteration of compositional
relations may be relevant to the constitution of an
existent. However, none of these transformations,
considered in isolation, establishes that the existent
has ceased to be the existent that it is.

To infer the loss of identity from compositional change
would require an independent account of the ontological
significance of the particular composition under
consideration.

Likewise, the preservation of composition cannot
establish numerical identity. Distinct existents may
possess equivalent or indistinguishable compositional
determinations.

The same compositional arrangement may therefore
be instantiated by numerically distinct existents.

Compositional invariance is consequently insufficient
to establish either the persistence or the loss of
ontological identity.

\subsection{The Limits of Compositional Explanation}

Composition provides an important dimension for
describing the constitution of an existent. Nevertheless,
it cannot be treated as an exhaustive explanation
of that existent.

An existent is not necessarily reducible to the sum
of its components, nor can its identity be established
solely through the preservation of their membership
or arrangement.

Such reduction would assume that every ontologically
relevant determination of an existent is contained
within its composition.

That assumption cannot be established through the
analysis of compositional invariance itself.

The ontological significance of composition must
therefore be acknowledged without granting it
explanatory exclusivity.

Compositional determinations may contribute to the
constitution and characterization of an existent,
but their preservation or transformation does not
independently settle the question of its identity.

\subsection{Synthesis}

The analysis establishes that compositional invariance
concerns the preservation of a specified compositional
determination across a defined transformation.

Its assessment requires distinguishing:

\begin{itemize}
    \item component membership from compositional organization;
    \item component persistence from compositional preservation;
    \item component replacement from the transformation
          of the entire composition;
    \item compositional invariance from structural invariance;
    \item compositional equivalence from numerical identity;
    \item compositional transformation from the loss
          of ontological identity.
\end{itemize}

The preservation of composition does not necessarily
establish the persistence of an existent, and the
transformation of composition does not necessarily
establish its loss of identity.

Compositional invariance must therefore remain a
determinate and limited category within the broader
analysis of invariance.

It identifies a possible mode of preservation in
the constitution of an existent, without reducing
the identity of that existent to the persistence
of its components or their compositional arrangement.

% =============================================================================
% End of Section 07
% =============================================================================

% =============================================================================
% SECTION 08 — INVARIANCE AND CONTINUITY
% =============================================================================

% =============================================================================
% Section 08 — Invariance and Continuity
% Book I — Ontology of Viability
% Chapter 02 — Invariance
%
% Purpose:
%   To establish the ontological distinctions between invariance,
%   continuity, and identity, and to examine their possible
%   relations under transformation.
%
% Conceptual Policy:
%   1. Invariance is not identical to continuity.
%   2. Continuity is not identical to numerical identity.
%   3. The absence of identified invariants does not establish
%      discontinuity or the loss of identity.
%   4. The preservation of invariants does not independently
%      establish continuity or numerical identity.
%   5. No universal criterion of continuity is assumed.
% =============================================================================

\section{Invariance and Continuity}
\label{sec:invariance-and-continuity}

\subsection{The Ontological Problem of Continuity}

Continuity constitutes a distinct ontological problem within
the analysis of transformation. It concerns the manner in which
an existent, or a determinate process involving an existent,
is connected across successive states, phases, or transformations.

The concept of continuity is frequently associated with the
preservation of certain determinations. When an existent changes
while some of its characteristics remain preserved, this
preservation may be interpreted as evidence of continuity.

However, such an interpretation requires further examination.

The preservation of a determination and the continuity of
an existent are not identical questions. A determination may
remain preserved across a transformation without establishing
that the existent itself has continued. Conversely, an existent
may undergo extensive transformation without the absence
of a particular preserved determination establishing a
discontinuity in its being.

Continuity must therefore be distinguished from invariance.

Invariance concerns the preservation of a specified
determination under a specified transformation. Continuity
concerns a relation of connectedness across successive
states or phases.

Although these concepts may be related, neither can be
reduced to the other.

\subsection{Continuity as a Relation Across Transformation}

Continuity is not simply the absence of change.

An existent may undergo continuous transformation while
its determinations change. The occurrence of transformation
does not necessarily interrupt continuity, just as the
absence of an immediately identifiable transformation does
not independently establish it.

The analysis of continuity therefore requires attention
to the relation between successive states or phases.

Such a relation may involve:

\begin{itemize}
    \item an ordered succession of states;
    \item a connection between successive phases;
    \item the persistence of a process through transformation;
    \item a relation of dependence between earlier and later
          states;
    \item a determinate organization of transformation.
\end{itemize}

These aspects represent possible dimensions of continuity.
They should not be treated as universally necessary conditions
for every form of continuity.

Different existents and processes may exhibit different
modes of connectedness. Consequently, continuity cannot
be defined through one undifferentiated criterion applicable
to every ontological domain.

The relevant relation must be specified according to
the existent or process under consideration.

\subsection{The Distinction Between Invariance and Continuity}

Invariance and continuity answer different analytical questions.

Invariance asks whether a specified determination remains
preserved under a specified transformation.

Continuity asks whether successive states or phases stand
in a determinate relation of connectedness.

The distinction becomes clear when their possible combinations
are considered.

A determination may remain invariant across successive
states, while the relation between the existents occupying
those states remains undetermined.

Likewise, an existent may undergo continuous transformation
even when none of its previously identified determinations
remains invariant.

Thus, invariance does not automatically establish continuity,
and continuity does not necessarily require the preservation
of a particular invariant.

The two concepts may be associated in particular cases,
but their relationship must be established rather than assumed.

\subsection{Continuity Without Identified Invariants}

The possibility of continuity without an identified invariant
is philosophically significant.

An existent may undergo a transformation in which its
previously identified determinations are altered. The absence
of a preserved determination does not, by itself, establish
that the existent has become discontinuous.

This conclusion follows from the distinction between
the preservation of determinations and the connectedness
of successive states.

A transformation may involve extensive alteration while
maintaining a determinate relation between successive
phases. The relevant continuity may therefore concern
the organization or connection of the transformation
rather than the preservation of a particular feature.

It is also possible that an invariant exists but has not
been identified. However, the ontological possibility
of continuity must not be made dependent upon the successful
identification of such an invariant.

The absence of an identified invariant is an epistemically
limited statement about the available determination of
preservation. It is not, by itself, an ontological proof
of discontinuity.

\subsection{Invariance Without Established Continuity}

The converse possibility must also be considered.

A specified determination may remain preserved across
different states or transformations without establishing
that the existents involved constitute one continuous
existent.

For example, distinct existents may possess the same
structural, compositional, or quantitative determination.
The preservation of that determination across them does
not establish a continuous passage from one existent
to another.

Similarly, an invariant may be preserved across a sequence
of states whose relation to one another has not been
ontologically established.

In such cases, invariance concerns the repeated preservation
of a determination, whereas continuity remains a separate
question concerning the connectedness of the states or
phases.

Consequently, the presence of invariants cannot serve
as an independent proof of continuity.

\subsection{Continuity and Numerical Identity}

Continuity must also be distinguished from numerical identity.

Numerical identity concerns whether an existent is one
and the same existent. Continuity concerns a relation
of connectedness across states or phases.

These concepts may be related, particularly when the
persistence of an existent through transformation is
under consideration. Nevertheless, continuity cannot
be treated as a sufficient definition of numerical identity.

A continuous process may involve distinct existents
at different stages. Likewise, a relation of continuity
between successive states does not, without further
ontological determination, establish that those states
belong to one numerically identical existent.

Conversely, the analysis of numerical identity cannot
be replaced by an analysis of continuity alone.

To infer numerical identity from continuity would require
an additional account of why the connected states belong
to one and the same existent rather than merely participating
in a continuous process or relation.

The distinction therefore remains necessary even where
continuity and identity appear closely associated.

\subsection{Continuity and Persistence}

Continuity is also distinct from persistence.

Persistence concerns the continued existence of an existent
across a temporal or transformational domain. Continuity
concerns the connectedness of successive states or phases.

Although persistence may involve continuity, the two
concepts are not interchangeable.

An existent may persist through transformations whose
successive phases are not adequately described by a simple
model of continuity. Conversely, a continuous process
may extend across a succession of distinct existents
without establishing the persistence of one numerically
identical existent.

The distinction prevents the assumption that every
continuous sequence constitutes the persistence of
one existent, or that every form of persistence must
be described through the same model of continuity.

The conditions under which persistence and continuity
coincide must therefore be examined separately.

\subsection{Continuity and Stability}

Continuity must not be confused with stability.

Stability concerns the maintenance of a condition,
configuration, or mode of organization under specified
circumstances or perturbations.

Continuity concerns connectedness across successive
states or phases.

A process may remain continuous while undergoing unstable
or rapidly changing configurations. Conversely, a stable
condition may be observed at distinct moments without
the stability itself establishing a continuous relation
between them.

Stability may contribute to the preservation of particular
determinations, but it is neither identical to continuity
nor a universal requirement for it.

The distinction is important because continuity should
not be interpreted as the persistence of an unchanged
condition.

\subsection{The Transformation of Continuity}

Continuity itself may be subject to different forms of
transformation.

The relation between successive states may remain preserved
while the determinations of those states change. It may
also undergo modification, interruption, or reorganization.

These possibilities require distinguishing transformation
within continuity from transformation of the relation
that constitutes continuity.

A change occurring within a continuous process does not
necessarily alter the continuity of that process. However,
a transformation affecting the relation between successive
phases may raise a distinct question concerning whether
the previous continuity remains applicable.

No universal conclusion can be derived merely from the
magnitude or number of changes involved.

The ontological significance of a transformation depends
on the particular relation of connectedness under
consideration.

\subsection{The Limits of Invariance as a Criterion of Continuity}

The preceding distinctions establish the limits of
using invariance as a criterion of continuity.

The preservation of an invariant may provide a relevant
determination when examining a continuous transformation.
Nevertheless, it cannot independently establish that
continuity exists.

Likewise, the failure to identify an invariant cannot
independently establish that continuity has ceased.

Any argument concerning continuity must therefore
identify the relation connecting the relevant states
or phases rather than relying exclusively on the
preservation or alteration of particular determinations.

This requirement prevents invariance from being assigned
an explanatory role beyond its proper analytical scope.

\subsection{Synthesis}

The analysis establishes the following distinctions:

\begin{itemize}
    \item Invariance concerns the preservation of a specified
          determination across transformation.
    \item Continuity concerns the connectedness of successive
          states or phases.
    \item Persistence concerns the continued existence of
          an existent across a specified domain.
    \item Numerical identity concerns whether an existent
          is one and the same existent.
    \item Stability concerns the maintenance of a condition
          or organization under specified circumstances.
\end{itemize}

These concepts may be related, but none can be substituted
for another without an additional ontological argument.

Continuity may occur without an identified invariant,
and invariance may be observed without established continuity.
Neither possibility determines numerical identity by itself.

The preservation of invariants is therefore one possible
aspect of the analysis of continuity, not its universal
foundation or sufficient criterion.

Continuity remains a distinct ontological problem whose
relation to identity, persistence, and transformation
must be examined without reducing it to the preservation
of unchanged determinations.

% =============================================================================
% End of Section 08
% =============================================================================

% =============================================================================
% SECTION 09 — FORMS AND DOMAINS OF INVARIANCE
% =============================================================================

% =============================================================================
% Project : THE_PHILOSOPHY_OF_VIABILITY
% Book    : Book I — Ontology of Viability
% Book ID : VS.B01
%
% File    : 2_MAINMATTER/PART_I_IDENTITY_AND_MANIFESTATION/
%           CHAPTER_02_INVARIANCE/sections/
%           section_09_forms_and_domains_of_invariance.tex
%
% Section : 09 — Forms and Domains of Invariance
% Chapter : 02 — Invariance
%
% Purpose :
%   To distinguish the analytical forms through which preservation
%   may be described from the domains in which invariance claims
%   are examined, and to clarify the limits of transferring
%   invariance claims between formal and ontological contexts.
%
% Conceptual Scope :
%   Forms of invariance, domains of invariance, formal domains,
%   ontological domains, four fundamental domains, three
%   intermediate hubs, cross-domain invariance, local and
%   cross-domain transformation, and the scope of invariance claims.
%
% =============================================================================

\section{Forms and Domains of Invariance}
\label{sec:forms-and-domains-of-invariance}

\subsection{The Need to Distinguish Forms from Domains}

The analysis of invariance requires the distinction between
two analytical dimensions: the forms through which preservation
is described and the domains within which a claim of preservation
is examined.

These dimensions address different questions. The forms of
invariance identify the kind of determination that is claimed
to remain preserved. The domains of invariance identify the
ontological or analytical field within which that preservation
is being considered.

The distinction is necessary because the description of what
remains preserved does not, by itself, determine the ontological
domain to which the determination belongs. Conversely, identifying
a domain does not establish which determination is preserved
within it or which form of invariance applies.

For example, structural invariance identifies a claim concerning
the preservation of an organizational determination. It does
not, by itself, establish whether the relevant structure belongs
to a formal system, an existent's structural determination,
or another context in which structural organization is examined.

A rigorous account must therefore specify both dimensions
without allowing one to substitute for the other. It must
also identify the transformation, respect, criterion, and
scope relevant to the particular claim.

\subsection{Forms of Invariance}

Forms of invariance identify analytically distinguishable
ways in which a determination may be claimed to remain
preserved under transformation. They concern the respect
in which preservation is asserted, not separate ontological
entities or necessarily independent kinds of being.

The following forms provide a working analytical vocabulary:

\begin{itemize}
    \item \textbf{Determinational Invariance:}
          The preservation of a specified determination
          or feature across a defined transformation.

    \item \textbf{Structural Invariance:}
          The preservation of a specified organizational
          pattern, arrangement, or structural constraint.

    \item \textbf{Relational Invariance:}
          The preservation of a specified relation between
          relata under a defined transformation.

    \item \textbf{Compositional Invariance:}
          The preservation of a specified determination
          concerning the composition of an existent.

    \item \textbf{Quantitative Invariance:}
          The preservation of a specified quantitative
          determination, such as magnitude, number, or ratio.

    \item \textbf{Qualitative Invariance:}
          The preservation of a specified qualitative
          determination under the transformation considered.

    \item \textbf{Configurational Invariance:}
          The preservation of a specified configuration
          or arrangement of determinations.

    \item \textbf{Relational-Structural Invariance:}
          The preservation of an organized pattern of
          relations, even where particular relata or
          individual relations may change, provided that
          the relevant organizational criterion remains
          satisfied.
\end{itemize}

These forms are analytical distinctions rather than an
exhaustive taxonomy. Their purpose is to clarify the respect
in which a preservation claim is made. They do not imply
that every existent exhibits every form, that every
transformation preserves at least one form, or that each
form has an independent ontological existence.

The classification also does not imply that the forms
are mutually exclusive. A single case may involve more
than one kind of preservation, depending on the determinations
and criteria under examination.

% =============================================================================
% Project : THE_PHILOSOPHY_OF_VIABILITY
% Book    : Book I — Ontology of Viability
% Book ID : VS.B01
%
% File    : 4_ASSETS/tables/
%           VS.B01.P01.C02.S09.SS02.P004.TL01.analytical_forms_of_invariance.tex
%
% Asset   : TL01 — Analytical Forms of Invariance
% Location: Chapter 02 / Section 09 / Subsection 02 / Paragraph 004
%
% Purpose :
%   To distinguish the principal analytical forms of invariance
%   by identifying their respective objects of preservation
%   and limits of inference.
%
% =============================================================================

\begin{table}[htbp]
    \centering
    \small
    \renewcommand{\arraystretch}{1.35}
    \setlength{\tabcolsep}{5pt}

    \begin{tabular}{|p{3.0cm}|p{5.2cm}|p{5.2cm}|}
        \hline
        \textbf{Form} &
        \textbf{Determination Examined} &
        \textbf{Limit of Inference}
        \\
        \hline

        Determinational &
        A specified determination or feature. &
        Does not establish the preservation of the
        existent as a whole or its numerical identity.
        \\
        \hline

        Structural &
        An organizational pattern, arrangement,
        or structural constraint. &
        Does not establish the preservation of
        every component or numerical identity.
        \\
        \hline

        Relational &
        A specified relation between relata. &
        Does not establish that every relatum
        remains unchanged or numerically identical.
        \\
        \hline

        Compositional &
        A specified determination concerning
        the composition of an existent. &
        Does not establish the preservation of
        every component or the whole existent's identity.
        \\
        \hline

        Quantitative &
        A specified magnitude, number, ratio,
        or other quantitative determination. &
        Does not establish qualitative, structural,
        or ontological preservation beyond its criterion.
        \\
        \hline

        Qualitative &
        A specified qualitative determination. &
        Does not establish quantitative or
        numerical identity.
        \\
        \hline

        Configurational &
        A specified configuration or arrangement
        of determinations. &
        Does not establish the preservation of
        every determination or component.
        \\
        \hline

        Relational-Structural &
        An organized pattern of relations under
        a specified organizational criterion. &
        Does not establish that every individual
        relation or relatum remains unchanged.
        \\
        \hline
    \end{tabular}

    \caption{Analytical forms of invariance and the respective
    determinations they examine. The listed forms are not
    mutually exclusive and do not constitute an exhaustive
    ontological taxonomy.}
    \label{tab:analytical-forms-of-invariance}
\end{table}

% =============================================================================
% END OF TABLE TL01
% =============================================================================

% =============================================================================
% Project : THE_PHILOSOPHY_OF_VIABILITY
% Book    : Book I — Ontology of Viability
% Book ID : VS.B01
%
% File    : 4_ASSETS/diagrams/
%           VS.B01.P01.C02.S09.SS02.P004.DG01.forms_and_domains_of_invariance.tex
%
% Asset   : DG01 — Forms and Domains of Invariance
% Location: Chapter 02 / Section 09 / Subsection 02 / Paragraph 004
%
% Purpose :
%   To visualize the analytical independence of forms of
%   invariance and domains of invariance.
%
% Conceptual Caution :
%   The figure represents possible analytical cross-classification.
%   It does not assert that every combination is instantiated
%   or ontologically established.
%
% Requires (preamble):
%   \usepackage{tikz}
%   \usetikzlibrary{calc, positioning, fit}
% =============================================================================

\begin{figure}[htbp]
    \centering

    \begin{tikzpicture}[
        font=\small,
        node distance=0.4cm,
        % --- base box ---
        box/.style={
            draw, rounded corners=3pt, align=center,
            inner sep=5pt, line width=0.6pt
        },
        % --- node roles (names avoid clashes with built-in TikZ keys) ---
        formbox/.style={box, text width=3.4cm, minimum height=2.3cm,
                        fill=blue!8, draw=blue!55!black},
        domainbox/.style={box, text width=2.4cm, minimum height=1.4cm,
                          fill=gray!10, draw=black!60},
        % --- frames, titles, labels ---
        vsgroup/.style={draw=black!45, dashed, rounded corners=5pt,
                        inner sep=8pt, line width=0.5pt},
        axistitle/.style={font=\small\bfseries, align=center},
        grouplabel/.style={font=\footnotesize\itshape, text=black!75,
                           align=center},
        % --- undirected link (no arrowheads: no priority between axes) ---
        vslink/.style={draw=black!50, dashed, line width=0.6pt}
    ]

    % ---------------------------------------------------------------
    % Form axis (left)
    % ---------------------------------------------------------------
    \node[formbox] (question) {
        What determination is claimed to remain preserved?
    };
    \node[formbox, below=of question] (examples) {
        Structural, relational, quantitative, qualitative,
        compositional, and other analytical forms
    };

    % ---------------------------------------------------------------
    % Domain axis (right): two columns of three, top-aligned with Form
    % ---------------------------------------------------------------
    \node[domainbox, anchor=north west] (d1) at
        ($(question.north east)+(1.8,0)$) {Formal Domain};
    \node[domainbox, below=of d1] (d2) {First Determination};
    \node[domainbox, below=of d2] (d3) {Structural Determination};

    \node[domainbox, right=of d1] (d4) {Relational Manifestation};
    \node[domainbox, below=of d4] (d5) {Material Manifestation};
    \node[domainbox, below=of d5] (d6) {Intermediate Hubs};

    % ---------------------------------------------------------------
    % Frames, titles, labels
    % ---------------------------------------------------------------
    \node[vsgroup, fit=(question)(examples)] (groupF) {};
    \node[vsgroup, fit=(d1)(d2)(d3)(d4)(d5)(d6)] (groupD) {};

    \node[axistitle, above=0.2cm of groupF] {FORM OF INVARIANCE};
    \node[axistitle, above=0.2cm of groupD] {DOMAIN OF EXAMINATION};

    \node[grouplabel, below=0.15cm of groupD] {
        Distinct domains of examination
    };

    % ---------------------------------------------------------------
    % Independence marker (symmetric, no arrowheads)
    % ---------------------------------------------------------------
    \node[circle, draw=black!60, fill=white, inner sep=2pt, font=\large]
        (times) at ($(groupF.east)!0.5!(groupD.west)$) {$\times$};
    \draw[vslink] (groupF.east) -- (times);
    \draw[vslink] (times) -- (groupD.west);

    % ---------------------------------------------------------------
    % Note (centred under both axes)
    % ---------------------------------------------------------------
    \coordinate (mid) at ($(groupF.south)!0.5!(groupD.south)$);
    \node[below=0.9cm of mid, text width=10.5cm, align=center,
          font=\footnotesize\itshape] {
        The form identifies the respect of preservation;
        the domain identifies the field in which that
        preservation is examined. Neither dimension
        determines the other.
    };

    \end{tikzpicture}

    \caption{Forms and domains constitute distinct analytical
    dimensions. The diagram presents their independence without
    asserting that every possible form-domain combination has
    been established.}
    \label{fig:forms-and-domains-of-invariance}
\end{figure}

% =============================================================================
% END OF DIAGRAM DG01
% =============================================================================

\subsection{The Non-Exclusivity of Forms}

The analytical forms of invariance may overlap because
a particular determination can be examined under more
than one respect.

A quantitative determination may be involved in a
structural arrangement. A relation may contribute to
a compositional organization. A configuration may
include qualitative and quantitative determinations
at the same time.

Such overlap does not make the distinctions redundant.
It indicates that different analytical descriptions
may address different aspects of one case of preservation.
The descriptions must nevertheless specify their respective
criteria so that the same feature is not counted several
times as though several independent invariants had
necessarily been established.

The forms should therefore not be interpreted as a
rigid classification in which every case must belong
to exactly one category. Nor should their analytical
distinction be taken to establish separate ontological
layers or independent entities.

Their usefulness depends on whether they clarify the
respect in which preservation is asserted. Where a
proposed category does not identify a distinct respect,
its status as a separate form requires reconsideration.

\subsection{Domains of Invariance}

A domain of invariance identifies the ontological or
analytical field within which a claim of preservation
is examined.

The domain does not specify, by itself, what is preserved.
It situates the determination under examination and
constrains the interpretation that may be assigned
to its preservation.

The same analytical form may be considered in different
domains. Structural invariance, for example, may be
examined within a formal system or in relation to the
structural determination of an existent. The use of
the same term does not establish that the two cases
have the same ontological status.

Likewise, a quantitative relation preserved within
a mathematical description does not, merely by being
formally valid, establish that the corresponding
determination is preserved in the existent being
described.

The domain must therefore be identified independently
of the form. Any transfer of a conclusion between
domains requires an account of the relation between
the determinations involved.

\subsection{Formal and Ontological Domains}

A distinction must be maintained between formal
domains and ontological domains.

A formal domain is constituted by the elements,
relations, operations, and rules specified within
a formal system. Mathematical, logical, and other
formally defined systems may provide domains in
which invariance can be demonstrated according
to explicit rules.

An ontological domain concerns determinations
belonging to an existent or to the constitution
of reality, independently of whether those
determinations are represented in a formal system.

Formal invariance may be established when a specified
property, relation, or structure is preserved under
a defined class of transformations within a formal
system. Such a result establishes a relation within
the scope of that system.

It does not automatically establish that an existent
instantiates the relevant formal structure, that
the formal transformation corresponds to an
ontological transformation, or that the property
is preserved in the existent itself.

The transition from a formal result to an ontological
claim therefore requires an additional account of
the correspondence between the formal elements
and the determinations of the existent. This account
must clarify what is represented, how the formal
transformation relates to the transformation under
consideration, and why the formal preservation is
relevant to the ontological claim.

This distinction does not diminish the value of
formal analysis. It specifies the scope of its
conclusions and prevents formal validity from
being treated as an independent proof of
ontological instantiation.

\subsection{Ontological Domains and Their Distinctness}

Within the proposed ontological articulation of
this book, four fundamental domains are distinguished:

\begin{enumerate}
    \item \textbf{First Determination}
    \item \textbf{Structural Determination}
    \item \textbf{Relational Manifestation}
    \item \textbf{Material Manifestation}
\end{enumerate}

These domains identify distinct dimensions of
ontological determination and manifestation.
Their distinction is intended to clarify different
aspects of the ontological articulation, not to
divide reality into isolated regions.

An existent may involve determinations relevant
to more than one domain. A transformation may
also affect determinations in several domains,
whether directly or through relations between them.

Nevertheless, the domains must not be treated
as interchangeable. A determination considered
within one domain cannot be assigned the same
ontological status as a determination considered
within another merely because similar language
is used to describe them.

The proposed distinctions also do not establish
that the four domains are exhaustive, mutually
exclusive in every possible application, or
arranged in a temporal or linear sequence.
Their precise relations remain subject to
the broader ontological analysis.

\subsection{The Four Fundamental Domains}

\subsubsection{First Determination}

First Determination concerns a foundational
dimension of determination within the proposed
ontological articulation.

An invariance claim situated in this domain
would concern the preservation of a determination
properly belonging to this dimension. Such a
claim must identify the relevant determination
and the respect in which it is said to remain
preserved.

It must not be inferred merely from the
preservation of a descriptive feature,
structural arrangement, relational pattern,
or material characteristic.

The precise nature and scope of any proposed
invariant within First Determination require
independent clarification. The identification
of this domain does not establish that a
particular invariant exists within it.

\subsubsection{Structural Determination}

Structural Determination concerns the structural
organization and determinate arrangement of
an existent.

Invariance within this domain may concern
the preservation of a specified organizational
pattern, structural relation, or constraint.

Such a claim must identify the relevant
structural respect and the transformation
under consideration. The preservation of
individual components does not automatically
establish the preservation of their organization.
Conversely, the preservation of an organizational
pattern does not necessarily establish that
every component remains unchanged.

Structural invariance must also remain distinct
from the identity of the existent. The preservation
of a structural determination does not independently
establish numerical identity, and structural
transformation does not independently establish
its loss.

\subsubsection{Relational Manifestation}

Relational Manifestation concerns the manifestation
of an existent through determinate relations.

An invariance claim within this domain may
concern the preservation of a specified relation,
relational determination, or organization of
relations.

The preservation of a relation does not
necessarily establish the preservation of
every relatum involved. Likewise, the alteration
of particular relations does not automatically
establish the disappearance of every relational
determination.

A claim must therefore specify the relata,
the relation under consideration, and the
respect in which preservation is asserted.
Where the organization of relations is involved,
it must be distinguished from the preservation
of each individual relation.

\subsubsection{Material Manifestation}

Material Manifestation concerns the material
manifestation of an existent.

An invariance claim within this domain may
concern specified material determinations,
quantities, configurations, or compositional
aspects.

Material invariance must not be equated with
immutability. A material existent may undergo
transformations while particular material
determinations remain preserved. Conversely,
the preservation of a material determination
does not establish that the existent itself
remains numerically identical.

The claim must identify the material determination,
the transformation, the criterion of preservation,
and the limits within which the preservation
is asserted. Material preservation cannot
be generalized beyond that scope without
additional justification.

\subsection{The Three Intermediate Hubs}

In addition to the four fundamental domains,
the proposed articulation distinguishes three
intermediate hubs:

\begin{enumerate}
    \item \textbf{Hub of Decree}
    \item \textbf{Hub of Signification}
    \item \textbf{Hub of Realization}
\end{enumerate}

These hubs are not presented as three additional
fundamental domains equivalent in status to
the four domains already identified. They
designate intermediate dimensions whose
ontological functions and relations require
their own analysis.

Their inclusion does not establish a linear
sequence connecting all seven designations.
Nor does it establish that each hub occupies
a fixed position between two particular
fundamental domains.

An invariance claim associated with a hub
must be formulated according to the
determination and function being examined.
The hubs must not be collapsed into the
four fundamental domains, and their
distinctive roles must not be erased by
treating all seven designations as equivalent
levels in a single hierarchy.

% =============================================================================
% Project : THE_PHILOSOPHY_OF_VIABILITY
% Book    : Book I — Ontology of Viability
% Book ID : VS.B01
%
% File    : 4_ASSETS/illustrations/
%           VS.B01.P01.C02.S09.SS08.P004.IL01.fundamental_domains_and_intermediate_hubs.tex
%
% Asset   : IL01 — Fundamental Domains and Intermediate Hubs
% Location: Chapter 02 / Section 09 / Subsection 08 / Paragraph 004
%
% Purpose :
%   To provide a visual orientation to the four fundamental
%   domains and three intermediate hubs in the proposed
%   ontological articulation.
%
% Conceptual Caution :
%   Spatial arrangement is illustrative only. It does not
%   establish hierarchy, temporal order, causal direction,
%   exhaustive classification, or fixed relations between
%   the domains and hubs.
%
% Requires (preamble):
%   \usepackage{tikz}
%   \usetikzlibrary{calc, positioning}
% =============================================================================

\begin{figure}[htbp]
    \centering

    % layout parameters (cm)
    \def\dx{4.0}   % horizontal offset of the fundamental domains from centre
    \def\dy{2.0}   % vertical offset of the top / bottom rows from centre
    \def\hx{1.8}   % horizontal offset of the two middle hubs from centre

    \begin{tikzpicture}[
        font=\small,
        node distance=0.4cm,
        % --- node roles ---
        fundamental/.style={
            draw=blue!55!black, fill=blue!8, rounded corners=5pt,
            line width=0.8pt, text width=3.2cm, minimum height=1.2cm,
            align=center, inner sep=6pt
        },
        hub/.style={
            draw=black!60, fill=gray!10, rounded corners=9pt, dashed,
            line width=0.8pt, text width=2.6cm, minimum height=1.2cm,
            align=center, inner sep=6pt
        },
        % --- non-directional association (no arrowheads) ---
        fieldline/.style={draw=black!55, densely dotted, line width=0.8pt},
        % --- legend swatches ---
        legfund/.style={draw=blue!55!black, fill=blue!8, rounded corners=2pt,
                        line width=0.6pt},
        leghub/.style={draw=black!60, fill=gray!10, rounded corners=4pt,
                       dashed, line width=0.6pt}
    ]

    % ---------------------------------------------------------------
    % Fundamental domains (four corners)
    % ---------------------------------------------------------------
    \node[fundamental] (first)      at (-\dx,  \dy) {\textbf{First Determination}};
    \node[fundamental] (structural) at ( \dx,  \dy) {\textbf{Structural Determination}};
    \node[fundamental] (relational) at (-\dx, -\dy) {\textbf{Relational Manifestation}};
    \node[fundamental] (material)   at ( \dx, -\dy) {\textbf{Material Manifestation}};

    % ---------------------------------------------------------------
    % Intermediate hubs (decree aligned with the top row)
    % ---------------------------------------------------------------
    \node[hub] (decree)        at (0,    \dy) {\textbf{Hub of Decree}};
    \node[hub] (signification) at (-\hx, 0)   {\textbf{Hub of Signification}};
    \node[hub] (realization)   at ( \hx, 0)   {\textbf{Hub of Realization}};

    % ---------------------------------------------------------------
    % Non-directional illustrative associations
    % ---------------------------------------------------------------
    \draw[fieldline] (first)      -- (decree);
    \draw[fieldline] (structural) -- (decree);
    \draw[fieldline] (first)      -- (signification);
    \draw[fieldline] (structural) -- (realization);
    \draw[fieldline] (relational) -- (signification);
    \draw[fieldline] (material)   -- (realization);
    \draw[fieldline] (relational) -- (realization);
    \draw[fieldline] (material)   -- (signification);

    % ---------------------------------------------------------------
    % Legend and note (centred under the bottom row)
    % ---------------------------------------------------------------
    \coordinate (mid) at ($(relational.south)!0.5!(material.south)$);

    \node[below=0.8cm of mid, font=\scriptsize] (legend) {
        \tikz[baseline=-0.5ex]\draw[legfund] (0,-0.12) rectangle (0.55,0.12);~Fundamental domain
        \quad\quad
        \tikz[baseline=-0.5ex]\draw[leghub] (0,-0.12) rectangle (0.55,0.12);~Intermediate hub
        \quad\quad
        \tikz[baseline=-0.5ex]\draw[fieldline] (0,0) -- (0.55,0);~Non-directional association
    };

    \node[below=0.45cm of legend, text width=10.5cm, align=center,
          font=\footnotesize\itshape] {
        Four fundamental domains and three intermediate hubs
        within one proposed ontological articulation.
        Their spatial arrangement does not establish a
        linear sequence or a hierarchy.
    };

    \end{tikzpicture}

    \caption{An illustrative orientation to the four fundamental
    domains and three intermediate hubs. The dotted associations
    are non-directional and do not specify fixed relations,
    causal connections, or ontological priority.}
    \label{fig:fundamental-domains-and-intermediate-hubs}
\end{figure}

% =============================================================================
% END OF ILLUSTRATION IL01
% =============================================================================

\subsubsection{Hub of Decree}

Hub of Decree identifies an intermediate
ontological dimension associated with decree.

Any proposed invariance within this hub
must concern a determination properly
belonging to its ontological function.
It cannot be assumed that preservation
within this dimension is equivalent to
the preservation of structural, relational,
or material determinations.

The specific criteria of invariance must
be established according to the nature
of the determination involved. The name
of the hub alone does not establish
which determinations are invariant or
how they are preserved.

\subsubsection{Hub of Signification}

Hub of Signification identifies an intermediate
ontological dimension associated with
signification.

It must not be reduced to language,
subjective meaning, mental representation,
or an act of interpretation. Its ontological
scope cannot be determined solely by
examining linguistic expressions or
psychological processes.

Consequently, invariance within this hub
cannot be defined exclusively through
the persistence of linguistic forms or
mental representations. Any claim must
address the determinations associated
with signification itself, while leaving
open the precise account of their
ontological constitution.

\subsubsection{Hub of Realization}

Hub of Realization identifies an intermediate
ontological dimension associated with
realization.

Invariance within this hub must be examined
according to the particular determinations
involved in realization. It should not
automatically be identified with material
persistence or with the preservation of
an already manifested configuration.

The relation between realization and
the other ontological dimensions requires
independent analysis. The designation
of this hub does not, by itself, establish
a universal mechanism or sequence through
which realization occurs.

\subsection{Cross-Domain Invariance}

An invariance claim may concern determinations
associated with more than one ontological domain.
Such a claim raises a question that must be
distinguished from the mere recurrence of
similar descriptions: are the same determinations
preserved across domains, or are different
determinations related through corresponding
forms of preservation?

These possibilities are not equivalent.

A determination that appears structurally
similar in two domains may have different
ontological functions, conditions, or criteria
of preservation. Similarity between descriptions
does not establish that one and the same
determination is preserved across the domains.

A cross-domain claim therefore requires:

\begin{itemize}
    \item identification of the determination
          considered in each domain;
    \item specification of the transformation
          relevant to each determination;
    \item clarification of the respect in which
          preservation is asserted;
    \item an account of the relation between
          the determinations across domains;
    \item justification for treating the result
          as cross-domain invariance rather than
          analogy or descriptive similarity.
\end{itemize}

Without such clarification, an apparent
cross-domain invariant may result from
similarity in the language of description
rather than an established ontological relation.

Cross-domain invariance is therefore a
claim requiring additional justification,
not a consequence of identifying comparable
features in different domains.

\subsection{Local and Cross-Domain Transformations}

Transformations may be examined within one
specified domain or in relation to determinations
extending across more than one domain.

A local transformation is one whose relevant
scope is restricted to a specified domain.
This does not imply that the transformation
is necessarily without consequences elsewhere;
it means that the claim under examination
is limited to the domain identified.

A cross-domain transformation involves
changes or dependencies that extend across
multiple domains. The relations through
which such changes are connected must
be specified rather than inferred from
their simultaneous occurrence.

This distinction is important because an
invariant may remain preserved within
one domain while relations between that
domain and others are altered.

Likewise, a transformation in one domain
may be associated with consequences in
another without producing an equivalent
transformation there.

Local preservation must therefore not
be confused with global preservation.
A claim of global invariance requires
an account of the relevant determinations
and transformations across every domain
included within its stated scope. It
cannot be inferred from the preservation
of one determination in one domain.

\subsection{The Scope of Domain-Specific Claims}

Every claim of invariance must have a
defined scope. At a minimum, the analysis
must identify the determination under
consideration, the form of invariance
claimed, the domain in which it is examined,
and the transformation and respect relevant
to the claim.

The scope may also require specification
of the criterion of preservation, the
conditions under which the claim applies,
and the limits beyond which the conclusion
is not intended to extend.

An invariant established within a limited
class of transformations cannot automatically
be generalized to all possible transformations.
Similarly, preservation within one domain
does not establish preservation across
the entire existent or across every domain
in which the existent may be considered.

A claim may be valid within a restricted
scope without being universally valid.
Such limitation does not necessarily
invalidate the claim; it determines
what the claim actually establishes.

The scope of validity must therefore
remain explicit throughout the analysis.
Where the scope changes, the conclusion
must be reassessed accordingly.

\subsection{The Limits of a Unified Classification}

The distinction between forms and domains
provides a framework for organizing
invariance claims. It does not imply
that every possible invariant can be
classified exhaustively within a closed
taxonomy.

Further distinctions among forms may
be required as additional modes of
preservation are examined. Likewise,
the analysis of domains may require
further clarification as their ontological
relations become more precise.

The forms listed in this section should
therefore be treated as an analytical
vocabulary rather than a completed
inventory. The four fundamental domains
and three intermediate hubs likewise
belong to a developing ontological
articulation, not an assertion that
every possible domain has already
been identified.

The framework is systematic in its
distinctions but remains open in its
scope. Its purpose is to prevent
conceptual confusion and unjustified
transfers, not to close the inquiry
into every possible form or domain
of invariance.

\subsection{Synthesis}

The analysis establishes two independent
but related dimensions for examining
invariance.

Forms of invariance identify the respect
or kind of determination claimed to
remain preserved. Domains of invariance
identify the ontological or analytical
field within which that preservation
is examined.

The distinction yields the following
conclusions:

\begin{itemize}
    \item A form of invariance does not
          determine its domain.

    \item A domain does not determine
          which form of invariance must
          occur within it.

    \item Formal invariance does not
          automatically establish
          ontological invariance.

    \item Local invariance does not
          establish global invariance.

    \item Invariance within one domain
          does not automatically transfer
          to another domain.

    \item The four fundamental domains
          and three intermediate hubs
          must not be treated as seven
          equivalent or linearly ordered
          ontological levels.

    \item No classification of invariance
          independently establishes the
          identity of an existent.
\end{itemize}

The analysis of invariance must consequently
preserve the distinction between what
is claimed to remain preserved and the
domain in which that preservation is
asserted. It must also distinguish the
formal validity of a preservation claim
from its ontological applicability.

Only through these distinctions can
invariance claims be evaluated without
confusing analytical form, domain,
formal representation, and ontological
determination. The classification
remains a means of clarifying the inquiry,
not a substitute for the arguments
required to establish particular
invariants.

% =============================================================================
% END OF SECTION 09 OF CHAPTER 02
% =============================================================================

% =============================================================================
% SECTION 10 — LIMITS OF INVARIANCE
% =============================================================================

% =============================================================================
% Section 10 — Limits of Invariance
% Book I — Ontology of Viability
% Chapter 02 — Invariance
%
% Purpose:
%   To establish the conceptual and ontological limits of invariance
%   and prevent its reduction to identity, persistence, continuity,
%   structure, or the exhaustive determination of an existent.
%
% Conceptual Policy:
%   1. Invariance is a limited mode of preservation.
%   2. Invariance cannot independently establish ontological identity.
%   3. The absence of identified invariants does not establish
%      the absence of identity.
%   4. Formal invariance does not automatically establish
%      ontological invariance.
%   5. No universal invariant is presumed to exist.
%   6. The limits of invariance do not diminish its analytical
%      and ontological significance.
% =============================================================================

\section{Limits of Invariance}
\label{sec:limits-of-invariance}

\subsection{The Need to Establish the Limits of Invariance}

The preceding analysis has established invariance as a
specific mode of preservation under transformation.
It has also distinguished invariance from identity,
difference, structure, composition, continuity, and
the domains within which preservation may be examined.

These distinctions make it necessary to establish the
limits of invariance explicitly.

Without such limits, invariance may be assigned an
explanatory role exceeding its proper scope. The preservation
of a determination may be treated as proof of the persistence
of an existent, while the transformation of that determination
may be interpreted as evidence of the loss of identity.

Both conclusions exceed what invariance alone can establish.

Invariance identifies the preservation of a specified
determination relative to a specified transformation.
It does not independently determine the complete ontological
status of the existent in which that determination occurs.

The purpose of this section is therefore not to diminish
the importance of invariance, but to establish the boundaries
within which its explanatory and analytical use remains valid.

\subsection{Invariance Is Not Identity}

The first and most fundamental limit concerns the distinction
between invariance and identity.

Identity concerns the existent as such. Invariance concerns
the preservation of a particular determination under
transformation.

The preservation of one or more determinations does not
establish that the existent possessing them is numerically
identical across the transformations under consideration.

Distinct existents may possess equivalent or even
indistinguishable invariant determinations. Their
structural, quantitative, qualitative, or compositional
features may coincide without making them one and the
same existent.

Conversely, an existent may undergo transformation in
which previously identified invariant determinations
are altered or cease to be preserved.

Such alteration does not independently establish that
the existent has lost its identity.

Consequently, neither the presence nor the absence of
an invariant is sufficient to determine numerical identity.

Identity must remain an independent ontological question.

\subsection{Invariance Is Not Persistence}

Invariance must also be distinguished from persistence.

Persistence concerns the continued existence of an existent
across a specified temporal or transformational domain.
Invariance concerns the preservation of a particular
determination within such a domain.

A determination may remain invariant across a sequence
of states without establishing that the same existent
persists throughout that sequence.

For example, distinct existents may exhibit the same
quantitative or structural determination. The preservation
of that determination across them does not establish
the persistence of one existent.

Likewise, an existent may persist through transformations
in which particular determinations are altered.

The persistence of an existent therefore cannot be reduced
to the preservation of an invariant.

Any claim concerning persistence requires an independent
account of the relation between the existent and its
successive states or transformations.

\subsection{Invariance Is Not Continuity}

Continuity concerns the connectedness of successive
states or phases. Invariance concerns the preservation
of a specified determination across transformation.

These concepts may coincide in particular cases, but
neither entails the other.

A process may remain continuous while its determinations
change. The absence of an identified invariant does not
necessarily establish a discontinuity.

Conversely, an invariant may be preserved across states
whose ontological connectedness has not been established.

The preservation of a determination therefore cannot
serve as a universal criterion of continuity.

Nor can the transformation or disappearance of an
invariant independently establish that continuity
has ceased.

Continuity must be examined according to the relation
connecting the relevant states or phases.

\subsection{Invariance Is Not Stability or Immutability}

Invariance is also distinct from stability and immutability.

Stability concerns the maintenance of a condition,
configuration, or organization under specified
circumstances or perturbations.

Immutability concerns the absence of change within
the relevant respect.

Invariance, by contrast, concerns preservation under
a specified transformation. Such preservation may
occur even when other determinations change.

An existent may therefore exhibit invariance without
being stable in every respect. It may also undergo
substantial transformation while preserving a particular
determination.

Likewise, the preservation of one invariant does not
establish that the existent is immutable.

The concepts must remain distinct because each addresses
a different relation between preservation and change.

\subsection{The Absence of a Universal Invariant}

No universal invariant should be presumed to exist
across all existents, transformations, and ontological
domains.

An invariant is always identified relative to a particular
determination, transformation, domain, and respect
of preservation.

A determination preserved in one case may not be preserved
in another. A transformation under which a particular
feature remains invariant may differ from a transformation
under which that feature changes.

Consequently, the existence of an invariant in one
existent or domain does not establish its universal
applicability.

Nor should the analysis assume that every existent
must possess at least one invariant under every possible
transformation.

Such a presumption would transform a limited analytical
concept into an unsupported universal ontological principle.

The possibility of invariance must be investigated
in each relevant case rather than imposed in advance.

\subsection{The Limits of Formal Invariance}

Formal invariance is established within a defined
formal system and relative to specified transformations
and preservation criteria.

Its validity may be mathematically or logically rigorous
within that system. However, formal validity alone
does not establish that the same invariant belongs
to an independently existing entity.

A formal representation may preserve a mathematical
relation while the ontological status of the corresponding
relation remains undetermined.

The transition from formal invariance to ontological
invariance therefore requires an account of the relation
between the formal system and the existent being described.

This requirement does not reject formal analysis.
It prevents the unqualified transfer of formal conclusions
into ontological claims.

Formal invariance may provide a precise description
of a possible pattern of preservation, but its ontological
significance must be independently established.

\subsection{The Limits of Observed Preservation}

The identification of an invariant may depend upon
the available descriptions, measurements, or analytical
methods through which a transformation is examined.

An observed determination may appear preserved within
a limited range of observations while its preservation
under a broader transformation domain remains
undetermined.

Likewise, a determination that appears to change under
one mode of description may remain preserved under
another appropriately specified criterion.

The distinction between observed preservation and
ontological preservation is therefore necessary.

Observed preservation provides evidence for an invariance
claim, but the scope of that claim cannot exceed the
conditions under which preservation has been established.

The absence of observed change is not automatically
proof of universal invariance.

Similarly, the failure to identify an invariant does
not establish its ontological absence.

These limitations concern the justification and scope
of invariance claims. They do not imply that invariance
is merely subjective or dependent upon observation.

\subsection{The Limits of Structural and Compositional Invariance}

Structural and compositional invariance provide
important descriptions of preservation, but neither
can serve as an exhaustive account of an existent.

Structural invariance concerns the preservation of
a specified organizational determination.

Compositional invariance concerns the preservation
of a specified determination relating to the constitution
of a composed existent.

Neither form establishes numerical identity by itself.

An existent may undergo structural reorganization
without the mere occurrence of that reorganization
establishing the loss of identity.

Likewise, the replacement or transformation of components
does not independently determine whether the existent
has ceased to be the existent that it is.

Conversely, the preservation of a structure or composition
does not establish that two descriptions or successive
states refer to one numerically identical existent.

The limits of these forms follow from the more general
distinction between the preservation of a determination
and the identity of the existent.

\subsection{The Limits of Cross-Domain Inference}

An invariant established within one domain cannot
automatically be transferred to another.

The preservation of a structural determination does
not establish the preservation of a corresponding
material determination.

Likewise, the preservation of a relational determination
does not automatically establish the preservation of
a determination belonging to First Determination.

Such transfers require an explicit account of the
relation between the relevant domains and determinations.

The same requirement applies to the four fundamental
domains and the three intermediate hubs.

Their distinction must be preserved when examining
invariance. The existence of a correspondence between
two domains does not establish that their determinations
are ontologically identical.

Cross-domain invariance must therefore be demonstrated
rather than presumed from similarity, analogy, or
formal correspondence.

\subsection{The Scope of Transformation}

An invariance claim is necessarily restricted by the
transformation domain within which preservation is
asserted.

A determination may remain invariant under a limited
class of transformations while changing under another.

The preservation of a determination under one transformation
therefore cannot be generalized to all possible
transformations without additional justification.

The same applies to the range of conditions under
which preservation is established.

An invariant may be conditional, domain-specific,
or restricted to a particular respect.

These restrictions do not invalidate the invariance
claim. They define its actual scope.

A precise account must therefore distinguish:

\begin{itemize}
    \item the transformation under which preservation
          is established;
    \item the transformations outside the established
          domain;
    \item the determination that remains preserved;
    \item the criterion by which preservation is assessed;
    \item the range within which the claim remains valid.
\end{itemize}

The absence of such specification makes an invariance
claim ambiguous and potentially misleading.

\subsection{Invariance Cannot Exhaust the Ontological Constitution of an Existent}

An existent cannot be exhaustively characterized
through the inventory of its invariant determinations.

Such an approach would assume that the ontological
constitution of an existent is reducible to those
determinations that remain preserved under selected
transformations.

This assumption is not justified.

An existent may possess determinations that change,
determinations that remain preserved, and determinations
whose relations to transformation have not yet been
established.

The ontological significance of an existent cannot
therefore be reduced to the set of features classified
as invariant.

Invariance identifies a particular relation between
determination and transformation. It does not provide
a complete account of being.

The analysis of invariance must consequently remain
situated within a broader ontological examination
of the existent.

\subsection{The Limits of Negative Inference}

The limits of invariance are especially important
when negative conclusions are drawn from the failure
to identify preserved determinations.

The absence of an identified invariant does not establish:

\begin{itemize}
    \item the absence of identity;
    \item the loss of numerical identity;
    \item the discontinuity of an existent;
    \item the termination of persistence;
    \item the absence of ontological organization;
    \item the absence of determinations not yet identified.
\end{itemize}

Such conclusions require independent arguments.

The failure to establish invariance may result from
the absence of preservation, the inadequacy of the
specified transformation domain, the selection of
an inappropriate criterion, or the limits of the
available analysis.

These possibilities must not be collapsed into one
ontological conclusion.

A negative invariance claim is therefore limited to
the determination, transformation, domain, and respect
that have actually been examined.

\subsection{The Limits of Positive Inference}

The same caution applies to positive conclusions.

The identification of an invariant does not establish:

\begin{itemize}
    \item the numerical identity of distinct states
          or existents;
    \item the persistence of one existent across every
          transformation;
    \item the continuity of a process;
    \item the immutability of the existent;
    \item the preservation of all other determinations;
    \item the identity of different ontological domains.
\end{itemize}

An invariant establishes preservation only within
its specified scope.

Any conclusion extending beyond that scope requires
additional ontological justification.

The presence of an invariant must therefore be treated
as a determinate finding rather than as a universal
explanation.

\subsection{The Positive Significance of Invariance}

Establishing the limits of invariance does not render
the concept philosophically marginal.

Invariance remains important because it allows the
analysis of preservation within transformation.

It makes possible the distinction between determinations
that remain preserved and those that undergo change.

It also provides a means of examining the conditions
under which particular features, relations, structures,
or compositional arrangements persist.

Without this distinction, the analysis of transformation
would risk treating all change as undifferentiated.

The proper limitation of invariance therefore strengthens
rather than weakens its analytical value.

By refusing to reduce identity, persistence, or continuity
to invariance, the ontological analysis preserves the
distinct significance of each concept.

\subsection{Synthesis}

The limits established in this section may be summarized
as follows:

\begin{itemize}
    \item Invariance is not identity.
    \item Invariance is not persistence.
    \item Invariance is not continuity.
    \item Invariance is not stability or immutability.
    \item No universal invariant is presumed.
    \item Formal invariance is not automatically ontological
          invariance.
    \item Observed preservation does not establish
          unrestricted ontological preservation.
    \item Invariance in one domain does not automatically
          transfer to another.
    \item The absence of an identified invariant does
          not establish the absence or loss of identity.
    \item The presence of an invariant does not establish
          numerical identity.
    \item Invariance cannot exhaust the ontological
          constitution of an existent.
\end{itemize}

Invariance must therefore be understood as a limited
but significant mode of preservation within transformation.

Its validity depends upon the specification of the
determination, transformation, domain, respect, and
criterion involved.

The recognition of these limits prevents invariance
from becoming an unjustified universal principle and
preserves the ontological distinction between the
preservation of determinations and the identity of
the existent.

% =============================================================================
% End of Section 10
% =============================================================================

% =============================================================================
% SECTION 11 — SYNTHESIS OF INVARIANCE
% =============================================================================

% =============================================================================
% Project : THE_PHILOSOPHY_OF_VIABILITY
% Book    : Book I — Ontology of Viability
% Book ID : VS.B01
%
% File    : 2_MAINMATTER/PART_I_IDENTITY_AND_MANIFESTATION/CHAPTER_02_INVARIANCE/sections/section_11_synthesis_of_invariance.tex
%
% Section : 11 — Synthesis of Invariance
% Chapter : 02 — Invariance
%
% Purpose :
%   To consolidate the ontological account of invariance, define
%   its proper conceptual and inferential scope, distinguish its
%   analytical dimensions, and identify the conclusions and
%   unresolved questions established by the chapter.
%
% Conceptual Scope :
%   Invariance; determination; transformation; preservation;
%   identity; numerical identity; persistence; continuity;
%   stability; immutability; forms and domains of invariance;
%   formal and ontological invariance; positive and negative
%   inference; unresolved ontological questions.
%
% =============================================================================

\section{Synthesis of Invariance}
\label{sec:synthesis-of-invariance}

\subsection{The Position of Invariance Within the Ontology of Identity}

The preceding sections have examined invariance as an
ontological problem concerning the preservation of
determinations under transformation. The analysis has
distinguished invariance from identity and has examined
its relations to determination, difference, structure,
composition, continuity, and the domains within which
preservation may be investigated.

The resulting account establishes invariance as a
significant but limited concept. It identifies a particular
relation between a determination and a transformation:
a specified determination remains preserved under a
specified transformation, within a specified scope.

This relation is significant because transformation does
not necessarily affect every determination in the same
way. Some determinations may remain preserved while
others are altered, replaced, reorganized, or left
undetermined by the analysis. Invariance makes this
differentiation conceptually explicit.

Nevertheless, the preservation of a determination does
not constitute the identity of the existent in which
that determination occurs. Nor does the alteration of
a determination independently establish the loss of
that identity.

Invariance must therefore be situated within the broader
ontology of identity without being made its foundation,
definition, or exhaustive criterion. Its contribution
concerns the analysis of preservation within transformation,
not the independent resolution of every question
concerning the existent.

\subsection{The Definition and Scope of Invariance}

On the basis of the preceding analysis, invariance is
defined as follows:

\begin{quote}
Invariance is the preservation of a specified
determination of an existent under a specified
transformation, relative to a defined domain,
respect, criterion, and range of validity.
\end{quote}

This definition expresses a relational and qualified
concept of preservation. It does not attribute
unrestricted sameness to an existent, nor does it
presume that every determination remains unchanged
whenever transformation occurs.

Several elements are essential to its interpretation.

First, invariance concerns a specified determination.
The determination must be distinguished from the
existent considered as a whole. The preservation of
one determination does not establish the preservation
of all the determinations belonging to that existent.

Second, invariance is asserted under a specified
transformation. A determination may remain preserved
under one transformation and be altered under another.
An invariance claim is therefore not meaningful as
an unrestricted attribution detached from transformation.

Third, the claim is made within a defined domain and
respect. These delimit the field and aspect in which
preservation is being examined. A claim concerning
one domain or respect cannot automatically be extended
to another.

Fourth, preservation must be assessed according to
a criterion. The criterion specifies what counts as
preservation rather than alteration, approximation,
resemblance, or another relation that may be relevant
to comparison.

Finally, the claim has a range of validity. Its
conclusion cannot exceed the transformations,
conditions, and determinations for which preservation
has been established.

These qualifications are not incidental additions
to an otherwise unrestricted concept. They determine
the content and scope of an invariance claim.

\subsection{Preservation and Identity}

The central ontological conclusion of this chapter
is that preservation and identity must remain
analytically distinct.

Identity concerns the existent as such. Invariance
concerns the preservation of a specified determination
under transformation. These concepts address different
questions, even when they are examined in relation
to the same existent.

Distinct existents may possess equivalent or
indistinguishable determinations. Their agreement
with respect to a specified invariant does not
establish their numerical identity.

Conversely, an existent may undergo transformations
in which previously identified invariant determinations
are altered or cease to be preserved. Such alteration
does not independently establish that the existent
has lost its identity.

The relation between identity and invariance must
therefore avoid two symmetrical reductions:

\begin{itemize}
    \item Identity cannot be reduced to the preservation
          of one or more determinations.
    \item The alteration or absence of an identified
          invariant cannot be treated as sufficient
          evidence for the absence or loss of identity.
\end{itemize}

This analytical distinction does not require identity
and invariance to be treated as unrelated in the
concrete existent. They may be examined together,
and findings concerning preservation may contribute
to arguments about identity. But the relevance of
such findings must be independently established.

The preservation of a determination is not identical
to the identity of that which possesses it.

\subsection{Transformation as a Condition of Invariance}

Invariance cannot be understood independently of
transformation. A preservation claim identifies what
remains preserved under a transformation; without
specifying the transformation, the scope of the
claim remains undetermined.

The same determination may be preserved under one
class of transformations and altered under another.
Consequently, invariance is transformation-relative
rather than an unqualified property attributed to
a determination in isolation.

This relativity does not make invariance subjective
or arbitrary. It means that the validity of a claim
must correspond to the transformation and conditions
under which preservation is asserted.

Transformation may also involve preservation and
alteration simultaneously. One determination may
remain preserved while another changes. A structural
relation may be maintained while components are
replaced; a compositional feature may be preserved
while the organization of the whole changes; or
a relational determination may remain preserved
while other aspects of the existent are modified.

Such possibilities must be assessed according to
the determinations and respects involved. Neither
preservation nor alteration should be assumed to
describe the existent exhaustively.

\subsection{The Analytical Requirements of an Invariance Claim}

The analysis of this chapter establishes that an
invariance claim must be sufficiently specified
before its validity and implications can be assessed.

At a minimum, the following elements must be
distinguished:

\begin{enumerate}
    \item \textbf{Subject of the Claim:}
          The existent, configuration, or analytical
          subject whose determination is being examined.

    \item \textbf{Specified Determination:}
          The particular determination alleged to remain
          preserved. It must not be confused with the
          existent considered as a whole.

    \item \textbf{Transformation:}
          The transformation under which preservation
          is asserted.

    \item \textbf{Transformation Domain:}
          The range or class of transformations included
          in the claim.

    \item \textbf{Domain and Respect:}
          The analytical or ontological domain and the
          particular aspect in which preservation is
          being examined.

    \item \textbf{Preservation Criterion:}
          The criterion by which preservation is
          distinguished from alteration or other
          relations of comparison.

    \item \textbf{Admissibility Conditions:}
          The grounds for treating the transformation,
          comparison, and criterion as applicable to
          the claim.

    \item \textbf{Ontological Grounding:}
          The justification for attributing ontological
          significance to the determination and the
          preservation relation under examination.

    \item \textbf{Scope of Validity:}
          The conditions and limits within which the
          conclusion is warranted.
\end{enumerate}

These elements are analytical requirements for
evaluating invariance claims. They are not independent
ontological entities, nor do they impose one mandatory
formal notation upon every investigation.

Their purpose is to make the content of a claim
determinate and to prevent conclusions from extending
beyond what its specification and justification
support.

\subsection{The Forms of Invariance}

The chapter has examined several analytically
distinguishable forms of invariance:

\begin{itemize}
    \item Determinational Invariance;
    \item Structural Invariance;
    \item Relational Invariance;
    \item Compositional Invariance;
    \item Quantitative Invariance;
    \item Qualitative Invariance;
    \item Configurational Invariance;
    \item Relational-Structural Invariance.
\end{itemize}

These forms identify different respects in which
preservation may be examined. They do not constitute
a closed inventory of independent ontological
categories.

The forms may overlap. A transformation may preserve
a quantitative determination while also preserving
a structural or relational determination. Conversely,
a determination may be preserved in one respect
while being altered in another.

The classification is therefore analytical rather
than exhaustive. Its categories organize questions
about what is preserved; they do not establish that
every instance of preservation must belong to one
and only one form.

The identification of a particular form also does
not establish the presence of other forms. Structural
invariance does not automatically establish
compositional invariance, and the preservation of
a quantitative determination does not establish
the preservation of qualitative or relational
determinations.

The classification remains open to refinement
where further analysis identifies additional
respects of preservation or requires more precise
distinctions among existing forms.

\subsection{The Domains of Invariance}

The forms of invariance must be distinguished from
the domains within which preservation is examined.

The ontological framework developed in this book
distinguishes four proposed fundamental domains:

\begin{enumerate}
    \item First Determination;
    \item Structural Determination;
    \item Relational Manifestation;
    \item Material Manifestation.
\end{enumerate}

It also distinguishes three intermediate hubs:

\begin{enumerate}
    \item Hub of Decree;
    \item Hub of Signification;
    \item Hub of Realization.
\end{enumerate}

These four domains and three hubs are not seven
equivalent levels arranged in one undifferentiated
sequence. Their conceptual functions and relations
must be considered according to the broader
ontological framework. Their enumeration does
not, by itself, establish a linear order, temporal
sequence, causal hierarchy, or exhaustive taxonomy.

Within the analysis of invariance, the distinction
between forms and domains is essential. Forms concern
the respect of preservation being examined; domains
concern the ontological or analytical field within
which the relevant determination is considered.

An invariance claim associated with one domain
or hub cannot automatically be transferred to
another. Such a transfer requires an explicit
account of the determinations involved and of
the relation that makes the proposed correspondence
ontologically relevant.

The same restriction applies to formal and
ontological domains. A formal result may provide
a precise description of preservation within a
mathematical or logical system, but it does not
automatically establish that the corresponding
determination is ontologically instantiated.

\subsection{The Relations Between Invariance and Neighboring Concepts}

The chapter has established relations between
invariance and several neighboring concepts without
identifying them with one another.

\begin{itemize}
    \item \textbf{Identity:}
          Invariance concerns the preservation of a
          specified determination; identity concerns
          the existent as such. Neither the presence
          nor the absence of an invariant independently
          settles numerical identity.

    \item \textbf{Persistence:}
          Invariance concerns preservation under
          transformation; persistence concerns the
          continued existence of an existent across
          a specified range of conditions. An invariant
          does not independently establish persistence.

    \item \textbf{Continuity:}
          Invariance concerns a preserved determination;
          continuity concerns a relation connecting
          successive states or phases. Neither entails
          the other.

    \item \textbf{Difference and Distinction:}
          Invariance identifies a determination that
          remains preserved under specified conditions.
          Difference and distinction concern other
          relations among determinations or existents.
          Their criteria must not be collapsed into
          a single notion of change.

    \item \textbf{Structure:}
          Structural invariance concerns the preservation
          of a specified organizational determination.
          It does not establish the identity of the
          existent or the preservation of every component.

    \item \textbf{Composition:}
          Compositional invariance concerns the preservation
          of a specified compositional determination.
          It does not necessarily establish the persistence
          of every component or the numerical identity
          of the composed existent.

    \item \textbf{Stability:}
          Stability concerns the maintenance of a condition,
          configuration, or organization under specified
          circumstances or perturbations. It is not
          identical to invariance.

    \item \textbf{Immutability:}
          Immutability concerns the absence of change
          within a specified respect, whereas invariance
          concerns preservation under a specified
          transformation. An invariance claim does not
          imply unrestricted immutability.
\end{itemize}

These distinctions preserve the independent analytical
roles of the concepts required for examining the
existent and its transformations. They also prevent
the analysis from assigning explanatory priority
to invariance merely because preservation can be
identified in a particular case.

\subsection{The Limits of Invariance}

The limits established in the preceding section
govern the interpretation of every invariance claim.

No universal invariant is presumed to exist across
all existents and transformations. The identification
of an invariant in one case does not establish its
universal applicability.

The preservation of an invariant does not
independently establish numerical identity,
persistence, continuity, stability, immutability,
or the preservation of all other determinations.

The absence of an identified invariant does not
independently establish the absence or loss of
identity, the termination of persistence, or the
discontinuity of an existent or process.

Formal invariance does not automatically establish
ontological invariance. Observed preservation does
not justify conclusions extending beyond the
conditions under which it has been established.
Likewise, preservation in one domain does not
automatically establish preservation in another.

Finally, the inventory of identified invariants
cannot be treated as an exhaustive account of the
ontological constitution of an existent.

These restrictions follow from the defined scope
of invariance. They do not diminish its significance;
they prevent its analytical contribution from being
mistaken for a complete ontological explanation.

\subsection{Consolidated Framework of Invariance}

The preceding analysis can be consolidated through
a framework that distinguishes the elements of
an invariance claim from the forms and domains
within which preservation is investigated.

The framework must preserve three analytical
dimensions:

\begin{enumerate}
    \item \textbf{Claim Specification:}
          What is alleged to be preserved, under which
          transformation, and according to which
          criterion and conditions?

    \item \textbf{Form of Preservation:}
          Which determination or respect of preservation
          is being examined?

    \item \textbf{Domain of Examination:}
          Within which analytical or ontological domain
          is the claim situated?
\end{enumerate}

These dimensions are related but not interchangeable.
The form of invariance does not determine the domain,
and the domain does not, by itself, determine the
form of preservation. Neither dimension replaces
the need to specify the transformation and criterion.

% =============================================================================
% Project : THE_PHILOSOPHY_OF_VIABILITY
% Book    : Book I — Ontology of Viability
% Book ID : VS.B01
%
% File    : 4_ASSETS/tables/VS.B01.P01.C02.S11.SS10.P004.TL01.consolidated_framework_of_invariance.tex
%
% Asset   : TL01 — Consolidated Framework of Invariance
% Type    : Table
% Location: Section 11, Subsection 10, Paragraph 004
%
% Purpose :
%   To consolidate the three analytical dimensions of an
%   invariance claim: claim specification, form of preservation,
%   and domain of examination.
%
% Dependencies :
%   Standard LaTeX table environments
%
% =============================================================================

\begin{table}[htbp]
    \centering
    \small
    \renewcommand{\arraystretch}{1.4}
    \setlength{\tabcolsep}{5pt}

    \caption{Three analytical dimensions of an invariance claim}
    \label{tab:consolidated-framework-of-invariance}

    \begin{tabular}{
        |p{0.20\textwidth}
        |p{0.34\textwidth}
        |p{0.34\textwidth}|
    }
        \hline

        \textbf{Dimension}
        &
        \textbf{Principal Elements}
        &
        \textbf{Analytical Question}
        \\ \hline

        Claim Specification
        &
        Subject; specified determination; transformation;
        transformation domain; respect; criterion;
        admissibility conditions; ontological grounding;
        scope of validity
        &
        What is claimed to be preserved, under which
        transformation, and within what limits?
        \\ \hline

        Form of Preservation
        &
        Determinational; structural; relational;
        compositional; quantitative; qualitative;
        configurational; relational-structural
        &
        What kind of determination or respect is
        preserved?
        \\ \hline

        Domain of Examination
        &
        First Determination; Structural Determination;
        Relational Manifestation; Material Manifestation;
        Hub of Decree; Hub of Signification;
        Hub of Realization
        &
        Within which analytical or ontological domain
        is preservation examined?
        \\ \hline

    \end{tabular}
\end{table}

% =============================================================================
% END OF TABLE TL01
% =============================================================================

% =============================================================================
% Project : THE_PHILOSOPHY_OF_VIABILITY
% Book    : Book I — Ontology of Viability
% Book ID : VS.B01
%
% File    : 4_ASSETS/diagrams/VS.B01.P01.C02.S11.SS10.P004.DG01.dimensions_of_an_invariance_claim.tex
%
% Asset   : DG01 — Dimensions of an Invariance Claim
% Type    : Diagram
% Location: Section 11, Subsection 10, Paragraph 004
%
% Purpose :
%   To represent the distinction and coordination of claim
%   specification, form of preservation, and domain of examination.
%
% Requires (preamble):
%   \usepackage{tikz}
%   \usetikzlibrary{arrows.meta, positioning, fit}
% =============================================================================

\begin{figure}[htbp]
    \centering
    \hyphenpenalty=10000 \exhyphenpenalty=10000\relax  % no hyphenation inside diagram labels

    \begin{tikzpicture}[
        font=\small,
        node distance=0.4cm,
        % --- base box ---
        box/.style={
            draw, rounded corners=3pt, align=center,
            inner sep=6pt, line width=0.7pt
        },
        % --- node roles (names avoid clashes with built-in TikZ keys) ---
        dimbox/.style={box, text width=3.3cm, minimum height=4.1cm,
                       fill=blue!8, draw=blue!55!black},
        assessbox/.style={box, text width=5.4cm, minimum height=1.1cm,
                          font=\small\bfseries,
                          fill=blue!16, draw=blue!55!black},
        resultbox/.style={box, text width=5.4cm, minimum height=1.1cm,
                          dashed, fill=green!8, draw=green!40!black},
        % --- frames and labels ---
        vsgroup/.style={draw=black!45, dashed, rounded corners=5pt,
                        inner sep=8pt, line width=0.5pt},
        grouplabel/.style={font=\footnotesize\itshape, text=black!75,
                           align=center},
        % --- arrows ---
        directed/.style={-{Stealth[length=2.5mm]}, line width=0.9pt}
    ]

    % ---------------------------------------------------------------
    % Three analytical dimensions (equal rank, no order among them)
    % ---------------------------------------------------------------
    \node[dimbox] (specification) {
        \textbf{Claim Specification}\\[5pt]
        \footnotesize
        Subject, determination, transformation,
        respect, criterion, conditions, scope
    };
    \node[dimbox, right=of specification] (form) {
        \textbf{Form of Preservation}\\[5pt]
        \footnotesize
        Determinational, structural, relational, compositional,
        quantitative, qualitative, configurational and combined forms
    };
    \node[dimbox, right=of form] (domaindim) {
        \textbf{Domain of Examination}\\[5pt]
        \footnotesize
        Four proposed fundamental domains
        and three intermediate hubs
    };

    \node[vsgroup, fit=(specification)(form)(domaindim)] (dims) {};
    \node[grouplabel, above=0.1cm of dims] {
        Distinct analytical dimensions, jointly coordinated
    };

    % ---------------------------------------------------------------
    % Assessment and result
    % ---------------------------------------------------------------
    \node[assessbox, below=0.9cm of dims] (assessment) {
        Invariance Assessment
    };
    \node[resultbox, below=0.9cm of assessment] (result) {
        Preservation claim within a specified scope
    };

    \draw[directed] (dims)       -- (assessment);
    \draw[directed] (assessment) -- (result);

    % ---------------------------------------------------------------
    % Note
    % ---------------------------------------------------------------
    \node[below=0.7cm of result, text width=10cm, align=center,
          font=\footnotesize\itshape] {
        The assessment does not independently establish
        numerical identity, persistence, continuity,
        or the total ontological constitution of the existent.
    };

    \end{tikzpicture}

    \caption{The three analytical dimensions of an invariance
    claim. They coordinate the assessment of preservation
    without becoming interchangeable or establishing conclusions
    beyond the specified scope.}
    \label{fig:dimensions-of-an-invariance-claim}
\end{figure}

% =============================================================================
% END OF DIAGRAM DG01
% =============================================================================

The table and diagram summarize the analytical
architecture developed in this chapter. They are
not intended as an exhaustive taxonomy of all
possible invariance relations. Their function is
to preserve the distinctions required for a
well-specified claim and to prevent conclusions
from being transferred between analytical dimensions
without justification.

\subsection{Established Conclusions}

The following conclusions have been established
through the analysis of this chapter:

\begin{enumerate}
    \item Invariance is the preservation of a specified
          determination under a specified transformation,
          relative to a defined domain, respect, criterion,
          and scope of validity.

    \item Invariance is analytically distinct from
          ontological identity.

    \item The preservation of an invariant does not
          independently establish numerical identity.

    \item The alteration or absence of an identified
          invariant does not independently establish
          the loss or absence of identity.

    \item Invariance is distinct from persistence,
          continuity, stability, and immutability.

    \item Transformation may involve preservation
          and alteration in different determinations
          or respects.

    \item Structural and compositional invariance
          concern distinct analytical respects of
          preservation.

    \item Forms of invariance and domains of invariance
          are distinct analytical dimensions.

    \item Formal invariance does not automatically
          establish ontological invariance.

    \item An invariance claim must be restricted to
          the transformations, conditions, and criteria
          for which preservation has been established.

    \item No universal invariant is presumed across
          all existents and transformations.

    \item The identified invariants of an existent
          do not exhaust its ontological constitution.
\end{enumerate}

These conclusions establish the conceptual position
of invariance within the ontology of identity. They
do not imply that every question concerning the
ontological significance of preservation has been
resolved.

\subsection{Open Questions and Further Investigation}

Several questions remain open and require further
ontological investigation.

First, the precise relation between invariance and
the determinations constitutive of an existent
requires further clarification. The present chapter
has established that invariance cannot replace
identity, but it has not determined whether some
forms of preservation may have a more fundamental
ontological significance than others.

Second, the conditions under which invariance may
be established across multiple ontological domains
remain unresolved. The distinction between the
four fundamental domains and three intermediate
hubs has been maintained, but the relations that
could warrant cross-domain preservation claims
require further examination.

Third, the relation between invariance and extensive
transformation remains open. The chapter has
distinguished preservation from identity and has
rejected the reduction of identity to invariant
determinations. It has not, however, provided a
complete account of the ontological conditions
under which an existent may undergo extensive
transformation.

Fourth, the possibility of multiple identity-related
determinations within one existent requires further
investigation. The identification of multiple
invariants does not establish multiple identities,
just as the identification of one invariant does
not establish that an existent has only one identity.

Fifth, the relation between formal invariance and
ontological invariance requires more precise
criteria for determining when a formal preservation
corresponds to an ontological determination.

Sixth, the possible existence of invariants with
a broad or cross-contextual range of applicability
remains open. The rejection of a presumed universal
invariant does not exclude the possibility that
particular invariants may prove to have extensive
scope.

These questions are not conclusions that should
be resolved by stipulation. They identify the
boundaries of what the present analysis establishes
and indicate where further conceptual and ontological
work is required.

\subsection{The Ontological Significance of Invariance}

The contribution of invariance to the ontology
of identity lies in clarifying how preservation
may occur within transformation without being
mistaken for identity itself.

It establishes a way to distinguish determinations
that remain preserved from those that are altered
under specified conditions. This distinction makes
it possible to examine transformation without
treating every change as an undifferentiated
alteration of the existent in its entirety.

An existent may exhibit preservation and alteration
at the same time, provided that the determinations
and respects involved are distinguished. The
preservation of one determination does not cancel
the alteration of another; nor does the alteration
of one determination eliminate the possibility
that others remain preserved.

% =============================================================================
% Project : THE_PHILOSOPHY_OF_VIABILITY
% Book    : Book I — Ontology of Viability
% Book ID : VS.B01
%
% File    : 4_ASSETS/illustrations/VS.B01.P01.C02.S11.SS13.P002.IL01.preservation_and_alteration_within_transformation.tex
%
% Asset   : IL01 — Preservation and Alteration Within Transformation
% Type    : Illustration
% Location: Section 11, Subsection 13, Paragraph 002
%
% Purpose :
%   To illustrate the coexistence of preservation, alteration,
%   and undetermined status among different determinations
%   during the transformation of one existent.
%
% Requires (preamble):
%   \usepackage{tikz}
%   \usetikzlibrary{arrows.meta, calc, positioning, fit, backgrounds}
% =============================================================================

\begin{figure}[htbp]
    \centering

    \begin{tikzpicture}[
        font=\small,
        node distance=0.25cm,
        % --- one determination (base) ---
        item/.style={
            draw=black!70, rounded corners=3pt, align=center,
            text width=3.0cm, minimum height=0.7cm,
            inner sep=3pt, line width=0.8pt
        },
        % --- status styles (differ by line type AND fill) ---
        neutral/.style={item, fill=white},
        preserved/.style={item, fill=green!12, draw=green!40!black,
                          line width=1pt},
        altered/.style={item, fill=blue!10, draw=blue!55!black, dashed},
        undetermined/.style={item, fill=white, dotted, line width=1pt},
        % --- frames, titles, labels ---
        configbox/.style={draw=black!45, fill=gray!5, rounded corners=6pt,
                          inner sep=8pt, line width=0.7pt},
        cfgtitle/.style={font=\small\bfseries, align=center},
        vsgroup/.style={draw=black!45, dashed, rounded corners=5pt,
                        inner sep=8pt, line width=0.5pt},
        % --- arrow ---
        directed/.style={-{Stealth[length=2.5mm]}, line width=0.9pt}
    ]

    % ---------------------------------------------------------------
    % Initial configuration
    % ---------------------------------------------------------------
    \node[preserved] (a1) {Determination A};
    \node[neutral, below=of a1] (b1) {Determination B};
    \node[neutral, below=of b1] (c1) {Determination C};

    % ---------------------------------------------------------------
    % Resulting configuration (same row positions as the initial one)
    % ---------------------------------------------------------------
    \node[preserved, right=3.8cm of a1] (a2) {Determination A};
    \node[altered, below=of a2] (b2) {Determination B$'$};
    \node[undetermined, below=of b2] (c2) {Determination C?};

    % ---------------------------------------------------------------
    % Frames behind the items
    % ---------------------------------------------------------------
    \begin{scope}[on background layer]
        \node[configbox, fit=(a1)(b1)(c1)] (initial) {};
        \node[configbox, fit=(a2)(b2)(c2)] (result) {};
    \end{scope}

    \node[cfgtitle, above=0.15cm of initial] {Initial Configuration};
    \node[cfgtitle, above=0.15cm of result]  {Resulting Configuration};

    % ---------------------------------------------------------------
    % Transformation
    % ---------------------------------------------------------------
    \draw[directed] (initial.east) -- (result.west)
        node[midway, above=2pt, font=\small\bfseries] {Transformation};

    % ---------------------------------------------------------------
    % Key (centred under both configurations, same styles as the items)
    % ---------------------------------------------------------------
    \coordinate (mid) at ($(initial.south)!0.5!(result.south)$);

    \node[altered, text width=2.4cm, font=\footnotesize,
          below=0.9cm of mid] (kalt) {Altered};
    \node[preserved, text width=2.4cm, font=\footnotesize,
          left=0.3cm of kalt] (kpres) {Preserved};
    \node[undetermined, text width=2.4cm, font=\footnotesize,
          right=0.3cm of kalt] (kund) {Undetermined};

    \node[vsgroup, fit=(kpres)(kalt)(kund)] (key) {};

    % ---------------------------------------------------------------
    % Note
    % ---------------------------------------------------------------
    \node[below=0.6cm of key, text width=10cm, align=center,
          font=\footnotesize\itshape] {
        One determination remains preserved, another changes,
        and the status of a third has not been established.
        These different outcomes do not by themselves settle
        the numerical identity of the existent.
    };

    \end{tikzpicture}

    \caption{Preservation and alteration may coexist within
    one transformation. The diagram distinguishes a preserved
    determination, an altered determination, and one whose
    status remains undetermined. It makes no independent claim
    about the identity of the existent.}
    \label{fig:preservation-and-alteration-within-transformation}
\end{figure}

% =============================================================================
% END OF ILLUSTRATION IL01
% =============================================================================

This differentiated analysis is significant for
the examination of structure, composition, relations,
and material determinations. It allows the
investigation to identify what remains preserved,
what changes, and what has not yet been sufficiently
determined.

The significance of invariance therefore lies
neither in explaining the whole existent nor in
providing a universal criterion of identity. It
lies in making particular relations of preservation
within transformation conceptually explicit.

A properly delimited concept of invariance is
more precise and more useful than an unrestricted
one. Its value depends on the adequacy of the
determination, transformation, criterion, and
scope specified in each claim.

\subsection{Concluding Statement}

The central conclusion of this chapter is that
invariance is a determinate mode of preservation
within transformation, not the foundation or
exhaustive criterion of ontological identity.

A determination may remain preserved while other
determinations of the same existent undergo change.
An existent may undergo transformation without
every determination being altered. Conversely,
the preservation of a particular determination
across distinct existents does not establish
their numerical identity.

Neither preservation nor alteration, considered
in isolation, settles the ontological question
of identity.

The analysis of invariance must therefore maintain
a precise distinction between the determination
that is preserved, the transformation under which
it is preserved, and the existent whose identity
is under consideration.

By establishing this distinction, Chapter 02
defines the proper conceptual and inferential
scope of invariance within the ontology of identity.
It also identifies the questions that remain open
without converting them into unsupported principles.

The chapter thus provides a basis for further
examination of identity and manifestation while
preserving the distinction between the identity
of an existent and the particular determinations
that may remain invariant through its transformations.

% =============================================================================
% END OF SECTION 11 OF CHAPTER 02
% =============================================================================

% =============================================================================
% CHAPTER CONCLUSION
% =============================================================================

\section*{Chapter Conclusion}
\addcontentsline{toc}{section}{Chapter Conclusion}

This chapter has examined invariance as a distinct ontological
problem concerning what remains unchanged under specified
transformations and in a specified respect.

The analysis has established that an invariance claim cannot
be adequately assessed without identifying its subject,
the transformations under consideration, the relevant
domain, the respect in which preservation is asserted,
and the criterion by which the claim is evaluated. These
conditions delimit the scope of the conclusion that may
legitimately be drawn.

The chapter has also distinguished invariance from identity,
difference, transformation, structure, composition, continuity,
persistence, stability, symmetry, conservation, regularity,
and immutability. Although these concepts may intersect,
their relations must be established rather than presumed.

In particular, the presence of an invariant does not by
itself prove the numerical identity of its subject. The
preservation of a structure, component, relation, or
determination cannot automatically be treated as proof
that the existent to which it belongs remains numerically
identical. Equally, the failure to identify an invariant
under a specified analysis does not establish that identity
is absent or that no invariant could be established under
another transformation, domain, or respect.

The distinction between formal and ontological invariance
also remains essential. Formal invariance is established
relative to a formal system and its operations. An ontological
claim requires an account of why the relevant relation
is applicable to being and not merely to the formal
representation through which it is expressed.

The conclusions of this chapter are therefore deliberately
limited. They establish conditions for the responsible
formulation and evaluation of invariance claims without
presupposing a universal invariant, a single underlying
ground, or an exhaustive account of the ontological role
of invariance.

The next chapter examines unity and plurality. Its inquiry
must not assume that invariance has already settled whether
an existent is one, whether its identity is one, or whether
a plurality of determinations, relations, structures, or
manifestations corresponds to a plurality of identities.
Those questions require independent analysis.

% =============================================================================
% END OF CHAPTER 02
% =============================================================================